\documentclass[11pt]{article}
\usepackage{fontspec}
\setmainfont{DejaVu Serif}
\setmonofont{DejaVu Sans Mono}
\usepackage{amsmath,amssymb}
\usepackage[margin=1in]{geometry}
\usepackage[hidelinks]{hyperref}
\usepackage{fvextra}
\setlength{\parindent}{0pt}
\setlength{\parskip}{0.6em}
\emergencystretch=4em
\title{Bounded geometric local products on excellent reduced germs\\Supported Cartier intersections and ordinary open locality}
\author{Independent whole-original P12.4 reconstruction}
\date{8 October 2026}
\begin{document}
\maketitle

\section*{Original target and reconstruction assignment}
The unchanged atomic target is P12.4 (T5): construct bounded mixed
Monge--Ampere products on singular and nonnormal normalized local boundaries,
with model independence, the correct intersection comparison, locality and
continuity. Compactness of the boundary alone supplies no such product.
The entire original family dossier and the augmented reconstruction assignment
are retained separately and literally in the accompanying task files.
The augmented request specifies proof audits; it changes no atomic target.
The sole incoming theorem is the entire qualified G331.

\section*{Theorem}
Let \(A\) be excellent local, \(R=A_{\rm red}\), and \(n=\dim R\).
Choose a proper primary \(\mathfrak a\) when \(n\ge1\). On the FULL
normalized centered semivaluation boundary, with every support retained,
there is a unique mixed positive Radon operator of arity \(p=n-1\)
on the bounded G331 geometric class \(D\), extending the independently
defined Cartier core and satisfying uniform weak and actual bounded
decreasing-net continuity. Decreasing-net limits are prescribed members of
\(D\). The operator is symmetric, additive and nonnegatively real homogeneous
in every nonpositive input. For actual negative-effective real curve-nef
primary-model Cartier cores its measure is
\[
\sum_Z \operatorname{ord}_Z(\mathfrak a)
 (D_1|Z\cdots D_p|Z\cdot Z)\,
 \delta_{\operatorname{ord}_Z/\operatorname{ord}_Z(\mathfrak a)}.
\]
Its pairing with every exceptional real Cartier test \(F\) on any common
primary refinement is the independently supported local degree
\(Q(D_1,\ldots,D_p,F)\). It depends only on the input functions.
For rational primary profiles the exact moments are
\[
\int v(J)\,d\mu(-v(I_1)/q_1,\ldots,-v(I_p)/q_p)
=\frac{e_R(I_1,\ldots,I_p,J)}{\prod q_i}.
\]
The reference tuple \((-1)^{[p]}\) has mass \(e_R(\mathfrak a)\).
The operator is local on every open set where corresponding bounded tuples
agree. Primary change \(T(v)=v/v(\mathfrak b)\), with transported inputs
\(u_{\mathfrak b}(Tv)=u_{\mathfrak a}(v)/v(\mathfrak b)\), obeys
\(\mu_{\mathfrak b}=T_*[v(\mathfrak b)\mu_{\mathfrak a}]\).
Finite normalization obeys closed residue-field weighted pushforward from
top-dimensional normal centers, with every conductor fiber retained and
lower-dimensional top-degree contribution zero. The curve operator has
zero inputs and is its actual Cartier reference measure; dimension zero
has empty boundary and zero measure, without negative arity.
The operator specializes exactly to published BFJ Section 4 at its smooth
affine CLOSED algebraically closed characteristic-zero hypotheses.

The theorem uses reduced top degree and G331's directed PRIMARY model system.
It supplies neither nilpotent generic lengths nor arbitrary-model cofinality,
general envelopes, energy, arbitrary finite-map calculus, or a literal
classical/CLD radial-extension identification. The mixed-zero obstruction
remains. This is an unverified complete original-solution candidate.

\section{n1. All-field primary mixed-multiplicity kernel on the reduced excellent germ; unused labels, same superficial quotients and closure invariance proved.}
\textit{Supports: G331.}

The only incoming theorem is the entire qualified G331, with its unchanged paper, annotation, ledger and scope. Let \(R=A_{\mathrm{red}}\) for excellent local \(A\), \(n=\dim R\ge1\), and let \(\mathfrak a\) be proper primary for the marked maximal ideal \(\mathfrak m\). Write \(X=X_{\mathfrak a}\) for the full compact Hausdorff section \(v(\mathfrak a)=1\). Every prime support is retained. G331 supplies the independently geometric continuous class \(C\), its equality with the uniform ideal closure, and its actual decreasing-net completion \(D\). All functions here are nonpositive; products take bounded \(D\) inputs. G331 grants no operator.

All multiplicities use degree \(n\) of the REDUCED ring. Modules of smaller dimension contribute zero. A positive-dimensional local ring has no zero-dimensional minimal component: its maximal ideal cannot be minimal while a smaller prime exists. Thus the zero-dimensional minimal-component exception in G331's valuative criterion does not arise here. Superficial quotients may be nonreduced, nonpure, and have embedded primes.

The ordinary published multigraded Hilbert theorem (complete current July 2024 Swanson--Huneke 17.4.2 and Definition 17.4.3) says that for finite \(M\) and proper primary \(I_1,\ldots,I_r\), the length \(\ell_R(M/I_1^{s_1}\cdots I_r^{s_r}M)\) is eventually a rational polynomial of degree \(\dim M\). Its degree-\(n\) form is
\[
H_M(s)=\sum_{|d|=n}\frac{e_R(I_1^{[d_1]},\ldots,I_r^{[d_r]};M)}{d_1!\cdots d_r!}s^d.
\]
With \(n\) separately labelled entries, \(e(I_1,\ldots,I_n)\) is the coefficient of \(s_1\cdots s_n\). Repeated equal entries recover ordinary multiplicity by substituting the sum of their variables. Unit labels mean independence of their variables and zero coefficients whenever those variables occur.

Unused labels need proof. Fix the unused primary \(B\)'s exponent at a sufficiently large \(c\). If \(K\) is a remaining primary label, choose \(N\) with \(K^N\subset B^c\). For the remaining product \(P_s\),
\[
P_sK^N\subset B^cP_s\subset P_s.
\]
The outside length polynomials differ only by a constant exponent shift and have the same top form. Squeezing on every positive rational ray proves the same top form for the middle polynomial. In the full polynomial, fixing the unused exponent leaves a degree-\(n\) part in the other variables coming exactly from terms with unused exponent zero. Polynomial identity proves deletion of unused labels, including repeated labels. With no remaining labels the degree-\(n>0\) coefficient is zero; units are immediate.

Relabelling proves symmetry. Substitution \(s\mapsto ls\) proves \(e(I^l,K_2,\ldots,K_n)=l e(I,K_2,\ldots,K_n)\). In the polynomial with \(n+1\) labels \(I,J,K_2,\ldots,K_n\), set the first two variables equal to \(s\). Only exponent vectors \((1,0,1,\ldots,1)\) and \((0,1,1,\ldots,1)\) contribute to \(s s_2\cdots s_n\). The unused-label result identifies their coefficients, giving
\[
e(IJ,K_2,\ldots,K_n)=e(I,K_2,\ldots,K_n)+e(J,K_2,\ldots,K_n).
\]
Iteration proves every power/product identity, with unit factors.

For integral-closure invariance suppose \(I\subset J\) and \(J\) is integral over \(I\). Homogeneous ideal equations make \(R[Jt]\) finite over \(R[It]\). Homogeneous module generators of degrees at most \(c\) give \(J^{m+c}\subset I^m\): in degree \(m+c\) their coefficients lie in \(I^{m+c-d}\subset I^m\). For \(m>c\) and any other primary product \(P_s\),
\[
I^mP_s\subset J^mP_s\subset I^{m-c}P_s.
\]
The outside length polynomials have identical top forms. Their difference on positive rational rays is lower order, so the middle form agrees on every such ray, hence identically. All mixed coefficients agree. Equal integral closures reduce to this case twice through \(I+J\). This is equality of forms, not a coefficient inequality deduced from a nonnegative polynomial; \((s-t)^2\) demonstrates why the latter inference fails.

For all-field monotonicity suppose \(J\subset I\) and keep \(K_2,\ldots,K_n\) fixed. Pass algebraically, if needed, to \(T=R[z]_{\mathfrak mR[z]}\), faithfully flat local with residue field \(k(z)\). Tensoring a finite-length composition series preserves length, with factors \(k(z)\). Dimension is preserved: an \(n\)-generator parameter ideal extends to an \(\mathfrak mT\)-primary ideal, giving \(\dim T\le n\) by height, while flat going-down lifts an \(n\)-chain below \(\mathfrak mT\), giving the reverse bound. Length polynomials and coefficients are therefore unchanged. No boundary base-change claim is made.

Over the infinite residue field, complete 17.2.2 provides complements of finitely many proper linear subspaces in \(K_n/\mathfrak mK_n\) for superficial choices for each of \((K_n,J,K_2,\ldots)\) and \((K_n,I,K_2,\ldots)\). Add avoidance of the current minimal primes. These added subspaces are proper: \(K_n\subset P+\mathfrak mK_n\) would imply \(K_n\subset P\) by Nakayama on \(K_n/(K_n\cap P)\), impossible for a primary ideal in positive dimension. Their finite intersection has a point over an infinite field. Choose one lift \(x\), using the SAME quotient for both families. Its dimension is \(d-1\): a quotient prime chain extends strictly down to a minimal prime avoiding \(x\), giving the upper bound; lifting a quotient parameter system with \(x\) gives the lower bound by height.

The colon calculation controls embedded primes. Multi-ideal Artin--Rees follows by finite homogeneous generation of the intersection module over the multi-Rees algebra. For superficial \(x\), its consequence (complete 17.2.4), using the common primary radical and Krull intersection, is
\[
(P_s:x)=(0:x)+P_{s-e_n},\qquad (0:x)\cap P_{s-e_n}=0
\]
for all large exponents. Artin--Rees first puts the colon in \((0:x)+P_{s-c}\); a sufficiently high \(K_n\)-power puts the remainder in the truncated range where superficiality applies. Intersecting the truncated superficial identities over arbitrarily large \(K_n\)-exponents and applying Krull intersection proves the second assertion. Thus \((0:x)\) injects into a finite primary quotient and has finite length. Multiplication by \(x\) gives
\[
\ell(R/P_s)-\ell(R/P_{s-e_n})
=\ell(R/(x,P_s))-\ell(0:x).
\]
For \(d>1\) the final constant does not affect the degree-\((d-1)\) coefficients. Reading the coefficient with one \(K_n\) entry removed gives the exact superficial reduction for BOTH families on the same quotient (complete 17.4.6). The unused-label result removes the zero \(K_n\) label.

Repeat in \(K_{n-1},\ldots,K_2\), always using simultaneous choices in the current common quotient. All images stay proper primary, and the arbitrary Noetherian quotient hypotheses remain valid. The final one-dimensional ring \(R'\) gives \(e_{R'}(JR')\) and \(e_{R'}(IR')\). The inequality between lengths of their powers, divided by the exponent, proves
\[
e_R(J,K_2,\ldots,K_n)\ge e_R(I,K_2,\ldots,K_n).
\]
For \(n=1\) this is direct. The same reduction and positive one-dimensional Samuel leading coefficient prove positive mixed multiplicities for proper entries. Units give zero. Symmetry proves every-slot monotonicity, and integral-closure invariance gives order for included integral closures. No regular-sequence or Cohen--Macaulay assumption was inserted.

We also need exact-sequence top-form additivity. For \(0\to M_1\to M_2\to M_3\to0\), multi-ideal Artin--Rees sandwiches \(P_sM_2\cap M_1\) between \(P_sM_1\) and \(P_{s-c}M_1\). The ray squeeze proves the same top form for its quotient lengths. The quotient exact sequence gives \(H_{M_2}=H_{M_1}+H_{M_3}\), with lower-dimensional modules contributing zero. This proves the application of complete 17.4.4 without asserting an injection after reduction modulo \(P_s\).

Finally G331's primary criterion on ALL \(X\) says \(v(J)\ge v(K)\) iff \(J\subset\overline K\), with units handled separately. Equal rational profiles give equal integral closures of the corresponding cross-powers. The laws above give equal mixed moments; profile inequalities give reversed ideal-moment order. These are mathematical proofs with the stated ordinary published Hilbert and superficial dependencies, not finite experiments.

\section{n2. Positive bounded moment measures on the whole G331 D domain, independent of every representation; actual prescribed-limit decreasing NET and uniform weak continuity.}
\textit{Supports: G331, n1.}

Put \(p=n-1\). For rational profiles \(u_i=-v(I_i)/q_i\), allowing unit ideals, define
\[
B_{\boldsymbol u}(v(J))=\frac{e_R(I_1,\ldots,I_p,J)}{q_1\cdots q_p}.
\]
For \(p=0\) this means \(e_R(J)\) with empty denominator one. N1 proves input-function independence, symmetry, addition, positive rational scaling, and ideal-test order.

Every rational ideal-difference test is \((v(J)-v(K))/d\). Equality with \((v(J')-v(K'))/d'\) gives
\[
v(J^{d'}(K')^d)=v(K^{d'}(J')^d)
\]
on ALL supports. Proper products have equal integral closures; equality with a unit forces both products to be units. The power/product laws prove equal proposed values, rational linearity, and all test-relation independence. A nonnegative difference gives closure inclusion and hence functional nonnegativity by n1; unit cases are immediate or impossible. Thus the whole rational lattice has a positive functional.

The constant \(1=v(\mathfrak a)\) belongs to that lattice. Positivity on \(r\pm f\), rational \(r>\|f\|_\infty\), gives \(|B(f)|\le\|f\|_\infty B(1)\) after decreasing \(r\). G331 supplies uniform density in real \(C(X)\). Norm-continuous extension is real-linear by rational scalar approximation. It is positive: for \(h\ge0\), approximate by lattice \(f_j\), add rational \(r_j>\|h-f_j\|\) tending to zero, and pass \(B(f_j)\ge-r_jB(1)\). Compact-Hausdorff Riesz representation gives a unique finite positive Radon measure \(\mu_{\boldsymbol u}\), of mass \(e_R(I_1,\ldots,I_p,\mathfrak a)/\prod q_i\). No metrizability is assumed.

Set \(E_J=e_R(\mathfrak a^{[p]},J)\), \(E_{\mathfrak a}=e_R(\mathfrak a)\). Comparing profiles in \([-M,0]\) to \(-M=-Mv(\mathfrak a)\), first rationally and then by rational upper limits, gives
\[
0\le B_{\boldsymbol u}(v(J))\le M^pE_J,\qquad
\mu_{\boldsymbol u}(X)\le M^p E_{\mathfrak a}.
\]
If two rational tuples have coordinate distances at most \(\epsilon\), use \(u_i-\epsilon\le w_i\) and the reverse inequalities. Expanding after adding the legitimate negative constant in every slot gives
\[
|B_{\boldsymbol u}(v(J))-B_{\boldsymbol w}(v(J))|
\le\Delta_p(M,\epsilon)E_J,\qquad
\Delta_p(M,\epsilon)=(M+\epsilon)^p-M^p. \tag{U}
\]
Each nonempty subset \(S\) of shifted slots contributes a nonnegative term at most \(\epsilon^{|S|}M^{p-|S|}E_J\); their sum is the stated polynomial. Rational upper errors remove restrictions on bounds. Set \(\Delta_0=0\), \(M^0=1\).

Approximate \(C\) tuples uniformly by rational primary profiles, using G331. A common finite bound and (U) make every ideal moment Cauchy and independent of approximation. Differences define the lattice limit. Positivity and mass bounds persist; Riesz gives the measure. Symmetry, addition and real nonnegative scaling pass to limits. If \(u\le w\), approximants within \(\eta\) satisfy \(u_{\rm app}-2\eta\le w_{\rm app}\); (U) removes this shift, proving order. The same bounds and (U) therefore hold throughout \(C\). Convergence of ideal moments and common mass bounds imply weak convergence for every continuous test by uniform lattice approximation. This proves uniform weak continuity for nets as well.

For bounded \(u_i\in D\), take ANY actual decreasing \(C\)-representations \(g_{i,\alpha_i}\downarrow u_i\). If \(u_i\ge-M\), all terms are already in \([-M,0]\). On the finite product index set, ideal moments increase and are bounded. Their suprema, and differences of these limits, define a positive norm-bounded lattice functional. Riesz gives its measure; density and mass bounds prove weak convergence of representing measures.

To prove independence, let \(f_{i,\beta_i}\downarrow u_i\) be another representation. Fix \(\boldsymbol\beta\) and \(\epsilon>0\). The directed open sets \(\{g_{i,\alpha_i}<f_{i,\beta_i}+\epsilon\}\) cover compact \(X\). Finite subcovers and directedness select one \(\alpha_i\) for each coordinate valid everywhere. Thus
\[
B_{\boldsymbol f_\beta}(v(J))
\le B_{\boldsymbol g_\alpha-\epsilon}(v(J))
\le\sup_{\boldsymbol\alpha}B_{\boldsymbol g_\alpha}(v(J))+\Delta_p(M,\epsilon)E_J.
\]
Decrease \(\epsilon\) to zero, take the supremum over \(\boldsymbol\beta\), and reverse the roles. Moments and measures coincide. The same cover argument for a fixed \(C\)-majorant \(h_i\ge u_i\) proves
\[
B_{\boldsymbol u}(v(J))
=\sup\{B_{\boldsymbol h}(v(J)):h_i\in C,\ h_i\ge u_i\}. \tag{M}
\]
Every such majorant belongs to \([-M,0]\); representing terms prove the reverse bound. This yields ideal-moment order on bounded \(D\): the majorants of a larger tuple form a smaller set. It asserts no order of measures.

Addition, scaling and symmetry on \(D\) follow by actual finite product \(C\)-representations and scalar moment limits. Sums of representing nets decrease to the required sums. Thus the negative shift expansion and (U) hold on bounded \(D\) as well. Uniform convergence INSIDE this established domain gives weak convergence. Increasing closure is not used.

For actual decreasing-net continuity let \(u_{i,\lambda}\in D\) decrease with common bound \([-M,0]\) to a SPECIFIED \(u_i\in D\). Order bounds the increasing moments above by the target. Fix a \(C\)-majorant tuple \(h\) and \(\epsilon>0\). The \(D\) functions are upper semicontinuous, so \(\{u_{i,\lambda}<h_i+\epsilon\}\) are directed open covers. Compactness selects a common \(\lambda\), giving
\[
B_{\boldsymbol h}(v(J))\le B_{\boldsymbol u_\lambda-\epsilon}(v(J))
\le B_{\boldsymbol u_\lambda}(v(J))+\Delta_p(M,\epsilon)E_J.
\]
Take the supremum over \(\lambda\), decrease \(\epsilon\) to zero, and take the supremum in (M). Every target ideal moment is recovered, and density and mass bounds prove weak convergence on every continuous test. Limit membership in \(D\) is a hypothesis. No nested-net flattening, arbitrary-infimum closure, or unbounded-energy limit is asserted.

\section{n3. Independent supported Cartier intersections on primary models: mixed-characteristic thickenings, actual Koszul symmetry, cycle projection and raw Rees equality.}
\textit{Supports: G331, n1.}

We define geometric intersections before ideal moments. On an actual primary blowup \(Y=\operatorname{Bl}_I\operatorname{Spec}R\), G331 supplies actual effective Cartier \(E_I\), ample \(\mathcal O_Y(-E_I)\), a reduced model whose components dominate the minimal components, and the exceptional Cartier order unit. Product-primary blowups give common refinements with actual invertible image ideals. Put \(p=n-1\).

For an effective exceptional integral Cartier \(E\), a power of \(\mathfrak m\) kills \(\mathcal O_E\): its coherent structure sheaf has closed-fibre support, so this follows on a finite affine cover. Its proper cohomology has finite \(R\)-length. Define
\[
\chi_R(E,\mathcal F)=\sum_j(-1)^j\ell_R H^j(E,\mathcal F).
\]
For line bundles \(L_1,\ldots,L_r\) on \(Y\), \(\chi_R(E,\mathcal O_E\otimes\bigotimes L_i^{t_i})\) is a polynomial of total degree at most \(p\). Indeed filter \(\mathcal O_E\) by \(\mathfrak m\)-powers. Its finitely many graded coherent quotients are on the proper \(k\)-scheme \(E_0\), of support dimension at most \(p\). Tensoring is exact; Euler additivity and complete Stacks 33.45.1 give the polynomial, with \(R\)-length equal to \(k\)-dimension on these quotients. This works in mixed characteristic without a coefficient field, and retains every nilpotent layer of \(E\).

For integral exceptional \(F_1,\ldots,F_p\) define \(Q_Y(F_1,\ldots,F_p,E)\) as the coefficient of \(t_1\cdots t_p\) in that polynomial with \(L_i=\mathcal O_Y(F_i)\). If \(p=0\), it is \(\ell_R H^0(E,\mathcal O_E)\): the supported proper zero-dimensional scheme is finite over its Artinian base and has no higher cohomology. Complete 33.45.2 applied to the filtration yields
\[
Q_Y(F_1,\ldots,F_p,E)
=\sum_{\dim Z=p}\ell_{\mathcal O_{Y,\eta_Z}}(\mathcal O_{E,\eta_Z})
 (F_1|Z\cdots F_p|Z\cdot Z). \tag{C}
\]
Here \(Z\) is integral reduced and proper over \(k\), and the parentheses are the numerical Euler-characteristic intersections of 33.45.3. At each top generic point, the filtration lengths add to the EXACT displayed local length: the residue field is unchanged by the \(\mathfrak m\)-quotient. Lower-dimensional filtration supports contribute zero. Thus all generic thickness weights are present. A fixed extra twist changes no degree-\(p\) coefficient.

Polynomial substitution proves additivity and symmetry of the first slots (33.45.5). The sequence
\[
0\longrightarrow\mathcal O_{E_1}(-E_2)
\longrightarrow\mathcal O_{E_1+E_2}
\longrightarrow\mathcal O_{E_2}\longrightarrow0
\]
and fixed-twist invariance prove additivity in the effective last slot. The Cartier order unit expresses every integral exceptional \(F\) as a difference of effective ones. This extends the last slot independently of decomposition, since equality of two differences is equality of their effective sums. Extend all slots to real Cartier divisors by multilinearity.

Full symmetry uses an ACTUAL supported Koszul complex. For effective exceptional \(E,F\), take \(K_E=[\mathcal O(-E)\to\mathcal O]\) and \(K_F=[\mathcal O(-F)\to\mathcal O]\), using their Cartier sections. Each is a locally free resolution of its quotient. Their tensor totals are identical up to the usual sign interchange. Tensor also with any line bundle \(L\). Compute its Euler characteristic on \(E\) and on \(F\); all terms and hypercohomology have finite-length cohomology, since support is proper over an Artinian base. The spectral-sequence Euler identity gives
\[
\chi_R(E,L)-\chi_R(E,L(-F))
=\chi_R(F,L)-\chi_R(F,L(-E)). \tag{K}
\]
Neither divisor must cut regularly on the other. Put
\(L=\mathcal O(F)^t\mathcal O(E)^s\bigotimes_{j=1}^{n-2}\mathcal O(G_j)^{r_j}\).
Set \(t=s=0\) and take the \(r_1\cdots r_{n-2}\) coefficient. The finite difference on the left extracts the \(t r_1\cdots r_{n-2}\) coefficient: total degree at most \(n-1\) excludes higher \(t\)-powers. The right extracts its \(s\) counterpart. Thus
\[
Q_Y(F,G_1,\ldots,G_{n-2},E)=Q_Y(E,G_1,\ldots,G_{n-2},F).
\]
Effective differences and real linearity give every transposition. For \(n=1\) there is no exchange. This pairing is supported locally, not a degree on global \(\mathrm{CH}_0(Y)\): for a DVR, the parameter divisor has \(Q=1\), although its global Chow class is zero.

For projection use \(\delta(P)=\dim R/P\). Catenarity and locality make this a dimension function; on finite-type models use its usual transcendence-degree extension. Excellence supplies the universally catenary hypotheses of complete Stacks 42.7.1. Top components have \(\delta\)-dimension \(n\), and closed-fibre dimension is ordinary \(k\)-dimension.

Let \(f:Y'\to Y\) be a primary refinement, finite normalization, or their normal-component lift. It is proper birational on each top component. Pullback of exceptional Cartier equations is regular, since these reduced components have generic points outside their support. At a dimension-\(p\) generic point of an exceptional divisor, only top components can occur: a smaller component dominates its own minimal component away from \(\mathfrak m\), and cannot contain a dimension-\(p\) closed-fibre subvariety. The local ring there is reduced one-dimensional. Embed such a local ring \(B\) into \(\bigoplus_{P\in\operatorname{Min}B}B/P\); the cokernel has finite length. Multiplication by a regular Cartier equation \(g\) has equal kernel and cokernel lengths on that finite-length module, giving
\[
\ell(B/gB)=\sum_P\ell((B/P)/g(B/P)).
\]
Thus the top divisor cycle is the sum of divisor cycles on the integral top components, whose generic reduced weights are one. Complete Stacks 42.26.3 applies to each component and gives the ACTUAL cycle equality
\[
f_*[f^*E]_p=[E]_p. \tag{P}
\]
Its function-field degree is one; residue-field degrees in cycle pushforward remain. Numerical projection, complete 33.45.7, on the proper closed-fibre components now gives
\[
Q_{Y'}(f^*F_1,\ldots,f^*F_n)=Q_Y(F_1,\ldots,F_n). \tag{PI}
\]
Components whose images have dimension below \(p\) give zero. Extend from effective last slots by differences. This proof never sends local degree through global Chow rational equivalence.

For the raw ideal diagonal let \(D_I=-E_I\). The standard finite Rees algebra has \(\mathcal O_Y(D_I)=\mathcal O_Y(1)\). Complete Stacks 30.14.3 gives \(H^0(Y,\mathcal O_Y(qD_I))=I^q\) for all large \(q\). Complete 30.16.1(4) gives higher cohomology vanishing. The actual Cartier sequence then yields
\[
\ell_R(I^{q-1}/I^q)=\chi_R(E_I,\mathcal O_{E_I}((q-1)D_I)).
\]
Its degree-\(p\) coefficients give \(Q_Y(D_I^{[p]},E_I)=e_R(I)\), or \(-Q_Y(D_I^{[n]})=e_R(I)\), including \(n=1\). No normality, CM property, residue-field infinitude or parameter-colength formula is used.

For \(I_1,\ldots,I_n\), take a common product blowup. The product ideal \(\prod I_j^{t_j}\) has divisor \(\sum t_jD_{I_j}\). Projection (PI) applies its diagonal formula on this common model. If \(H\) is the degree-\(n\) length form, \(e(\prod I_j^{t_j})=n!H(t)\), as seen by substituting exponents \(q t_j\). Equating polynomials on positive integer \(t_j\), then reading \(t_1\cdots t_n\), proves
\[
-Q_Y(D_{I_1},\ldots,D_{I_n})=e_R(I_1,\ldots,I_n). \tag{I}
\]
This is equality of whole polynomials; it makes no coefficientwise positivity inference.

\section{n4. Independent positive geometric atom measures and refinement invariance, then full real-core/ideal moment comparison by an ACTUAL fixed-model image identity.}
\textit{Supports: G331, n1, n2, n3.}

We next construct the core measure independently. Normalize a common primary model \(Y\), with finite normalization \(N\to Y\). Use only its top normal components; lower components and their boundary points remain in the domain but have zero degree-\(n\) contribution. Let \(Z\) run through their exceptional prime divisors of dimension \(p\). These are finitely many integral proper \(k\)-schemes, with DVR generic rings. Put \(b_Z=\operatorname{ord}_Z(\mathfrak a)>0\), \(\nu_Z=\operatorname{ord}_Z/b_Z\in X\). The valuation is pulled back from the fraction field of its minimal component, with that component as support. For negative-effective real curve-nef core divisors \(D_1,\ldots,D_p\), define
\[
\mu_Y^{\rm geom}(D_1,\ldots,D_p)
=\sum_Z b_Z(D_1|Z\cdots D_p|Z\cdot Z)\delta_{\nu_Z}. \tag{A}
\]
Pullbacks to \(N\) are understood. For \(p=0\), the intersection is \(\deg_k Z=[\kappa(Z):k]\).

Positivity precedes ideal comparison. G331's qualified real nef-plus-ample and finite-span rationalization give negative rational ample divisors approaching each \(D_i\) in coefficient space, by subtracting small \(E_I\)-multiples and rounding negative-effective coefficients upward. Finite normalization pulls ample bundles to ample bundles; restriction to \(Z\) stays ample. Clearing rational denominators and complete 33.45.9 give positive mixed numerical intersections. Their finite coefficient polynomials converge, so the weights in (A) are nonnegative. This argument also covers arbitrary real nef entries and needs no definition of ideal positivity.

For an exceptional real Cartier \(F\) on this model, G331's evaluation at \(\nu_Z\) is \(\operatorname{ord}_Z(F)/b_Z\). Normal divisor cycles and n3's projection yield the geometric identity
\[
\int_X H_F\,d\mu_Y^{\rm geom}=Q_Y(D_1,\ldots,D_p,F). \tag{G}
\]
Here \(H_E=v(I)\) for an effective \(E_I\); negative divisors evaluate negatively.

The atomic measure is invariant under refinements independently of algebraic moments. If \(Y'\to Y\) is a primary refinement, its normalization \(N'\) maps to \(N\): on each dominant normal component, integral functions defining \(N\) belong to the integrally closed local rings of \(N'\). This lift is proper (use the closed graph over \(Y\)) and birational on top components. For a prime \(Z'\) mapping dominantly to a dimension-\(p\) prime \(Z\) of \(N\), its DVR dominates \(\mathcal O_{N,\eta_Z}\) in the SAME fraction field. They are equal: every field element is a unit times a power of the lower uniformizer, and its sign of order is preserved. Their orders, residue fields, \(b\)'s and valuations therefore coincide. Apply cycle projection (P) to the pulled-back effective primary \(E_I\), whose order at \(Z\) is positive. It forces exactly one such prime over each \(Z\), with degree one. Numerical projection gives its mixed weight equal to the old weight. Every other \(Z'\) has image dimension below \(p\); its \(p\) pulled-back line bundles have intersection zero by 33.45.7. Thus (A) is invariant as a measure on \(X\). The same proof handles common refinements and auxiliary normal lifts, with every top component considered separately. In particular (G) can now be tested against Cartier functions determined on ANY common refinement.

We prove the bridge to n2's algebraic measures using actual image ideals, not only valuation identities. Fix \(Y=\operatorname{Bl}_I R\). Qualified G331 gives \(T=H^0(Y,\mathcal O_Y)\) finite birational, \(T/R\) supported at \(\mathfrak m\). Choose \(h\) with \(\mathfrak m^h(T/R)=0\) and fix ONCE
\(\mathfrak c_0=\mathfrak m^{h+1}\), so \(\mathfrak c_0T\subset\mathfrak m\) as elements of \(R\). For a negative rational ample exceptional \(L\) and sufficiently large divisible \(q\), \(\mathcal O_Y(qL)\) is globally generated. Set
\[
F_q=H^0(Y,\mathcal O_Y(qL))\subset T,\qquad J_q=\mathfrak c_0F_q\subset R.
\]
This is an ideal, proper because \(\mathfrak c_0T\subset\mathfrak m\), and unit off \(\mathfrak m\) by localization of global sections and triviality of the divisor there. Hence it is proper primary.

Fix \(Y'=\operatorname{Bl}_{I\mathfrak c_0}R\), dominating \(Y\) and \(\operatorname{Bl}_{\mathfrak c_0}R\); write \(D_0\) for the negative divisor of the latter. Global generation means equality of the IMAGE subsheaf of \(F_q\otimes\mathcal O_Y\) in \(\mathcal O_Y\) with the invertible \(\mathcal O_Y(qL)\). Pullback preserves its surjection. The pullback inclusion into \(\mathcal O_{Y'}\) remains injective, since its Cartier equation stays regular on every reduced dominating component. Multiplying by the invertible image of \(\mathfrak c_0\) proves the exact ideal identity
\[
J_q\mathcal O_{Y'}=\mathcal O_{Y'}(qf^*L+D_0). \tag{S}
\]
Consequently the ACTUAL blowup universal property gives \(Y'\to\operatorname{Bl}_{J_q}R\). This is stronger than equality of valuation functions.

For simultaneous \(L_1,\ldots,L_n\) and exponents \(q_i\), n3's (I) and projection give
\[
\frac{e(J_{1,q_1},\ldots,J_{n,q_n})}{q_1\cdots q_n}
=-Q_{Y'}(f^*L_1+D_0/q_1,\ldots,f^*L_n+D_0/q_n). \tag{S2}
\]
On every support, (S) also gives
\[
-v(J_q)/q=H_L-v(\mathfrak c_0)/q. \tag{S3}
\]
The error is uniform because \(\mathfrak c_0\) is PRIMARY. The right side of (S2) converges by the finite-dimensional Cartier polynomial on ONE fixed model.

For real negative curve-nef \(D_i\), take negative rational ample \(L_{i,j}\to D_i\) in their fixed finite coefficient spans, as in the positivity argument, and choose generated divisible \(q_{i,j}\to\infty\). The correction and refinement are unchanged. Both (S2)'s Cartier polynomials and (S3)'s function profiles converge. N2's uniform moment continuity, including its mass bound for varying continuous tests, gives
\[
\int_X H_{D_n}\,d\mu^{\rm alg}(H_{D_1},\ldots,H_{D_p})
=Q_Y(D_1,\ldots,D_p,D_n). \tag{B}
\]
Thus no uniform-function-to-intersection convergence was assumed.

For arbitrary integral exceptional \(F\), choose \(N\) so \(F-NE_I\) and \(-NE_I\) are both negative ample, using G331's Cartier order unit and ample twisting. Their difference is \(F\); linearity extends (B) to this test and then to every real exceptional Cartier test. A primary ideal test is determined on a common primary refinement. The independent atomic invariance proved above, (G), and (B) identify all its moments on that refinement. G331's dense ideal-difference lattice therefore gives
\[
\mu_Y^{\rm geom}(D_1,\ldots,D_p)
=\mu^{\rm alg}(H_{D_1},\ldots,H_{D_p}). \tag{COMPARE}
\]
This proves independence from all core presentations giving the same functions, as well as exact geometric Cartier comparison and primary-model independence. The algebraic extension of n2 is consequently the bounded product extending this independently constructed geometric core.

\section{n5. Full bounded ordinary open-set locality on all supports, from actual ratio subideals and two compact-cover limits.}
\textit{Supports: G331, n2, n4.}

Assume \(p\ge1\); zero arity is automatic. For rational profiles synchronize denominators and write \(u=-v(A)/q\), \(w=-v(B)/q\). Their maximum is \(-v(A+B)/q\). Take a common primary blowup making \(A,B,A+B\), the reference primary, and every other input ideal invertible, then normalize. Unit factors are harmless; the extra reference factor ensures an actual primary model even if other ideals are units. On this normal model put \(S=(A+B)\mathcal O_N\) and define ACTUAL ratio subideals
\[
L=A\mathcal O_N/S\subset\mathcal O_N,\qquad
M=B\mathcal O_N/S\subset\mathcal O_N,\qquad L+M=\mathcal O_N.
\]
They are invertible ideals, since all three original images are invertible. At an exceptional prime \(Z\) with \(u(\nu_Z)>w(\nu_Z)\), one has \(\operatorname{ord}_Z(A)<\operatorname{ord}_Z(B)\), so \(\operatorname{ord}_Z(M)>0\). Every local generator of \(M\) restricts to zero at the generic point of integral reduced \(Z\), hence to zero on its entire open restriction to \(Z\). Thus the entire map \(M|Z\to\mathcal O_Z\) is zero. Restricting \(L+M=\mathcal O_N\) gives a surjection \(L|Z\to\mathcal O_Z\), necessarily an isomorphism of line bundles. Therefore
\(\mathcal O_N(D_A-D_{A+B})|Z\) is trivial. In n4's independently geometric atomic formula, every mixed weight at an atom in the strict set is unchanged upon replacing the selected input \(u\) by \(\max(u,w)\). On the same common finite divisorial set, this proves
\[
\mu(\max(u,w),u_2,\ldots,u_p)|_{\{u>w\}}
=\mu(u,u_2,\ldots,u_p)|_{\{u>w\}}. \tag{MAX}
\]
This uses the actual subideals and their images, not just Weil coefficients or moment order.

For continuous \(C\) inputs approximate all entries uniformly by rational profiles. If a continuous test \(h\) has compact support \(K\) contained in the open set \(\{u>w\}\), the gap has a positive minimum on \(K\). Accurate approximants have the same strict inequality throughout \(K\). Apply (MAX) and n2's uniform weak continuity to pass the integrals against \(h\) to the limit. Thus (MAX) holds for \(C\) inputs. Now keep selected \(u,w\in C\) fixed and represent the other bounded \(D\) slots by actual decreasing \(C\)-nets. The strict open set and compactly supported test are fixed; n2's decreasing-net theorem proves (MAX) with all other slots in bounded \(D\).

Let bounded \(u,w\in D\) agree on an open \(U\), and fix the other bounded \(D\) slots. Choose actual decreasing \(C\)-representations \(u_\alpha\downarrow u\), \(w_\beta\downarrow w\), all with a common bound. Fix a continuous \(h\) with compact support \(K\subset U\), and \(\epsilon>0\). For each fixed \(\alpha\), the open sets
\[
\{w_\beta-\epsilon<u_\alpha\}
\]
increase and cover \(K\), because their pointwise limits there satisfy \(w-\epsilon=u-\epsilon<u_\alpha\). Compactness selects one \(\beta_0\). Its strict set is an open neighborhood \(V\) of \(K\), fixed for this \(\alpha,\epsilon\), and it remains strict for every \(\beta\ge\beta_0\). The continuous-selected-slot version of (MAX) gives
\[
\int h\,d\mu(\max(u_\alpha,w_\beta-\epsilon),\ldots)
=\int h\,d\mu(u_\alpha,\ldots).
\]
First decrease \(\beta\). The selected functions decrease to \(\max(u_\alpha,w-\epsilon)\), which belongs to bounded \(D\) by the actual representation or finite-max stability. N2's theorem applies with this prescribed limit. Then decrease \(\alpha\); the limit is the known bounded \(D\) function \(\max(u,w-\epsilon)\). Applying the same theorem again gives
\[
\int h\,d\mu(\max(u,w-\epsilon),\ldots)=\int h\,d\mu(u,\ldots). \tag{L1}
\]
These are two justified limits, not a flattening of nested nets. Interchange \(u,w\) to obtain the analogous identity for \(\max(w,u-\epsilon)\).

Both maximum functions are within \(\epsilon\) uniformly of the already known \(\max(u,w)\in D\). Although they increase as \(\epsilon\) decreases, n2's uniform estimate (U) applies INSIDE the established bounded domain. Letting \(\epsilon\) tend to zero in the two identities gives equality of the integrals for \(\mu(u,\ldots)\) and \(\mu(w,\ldots)\).

Such tests determine Radon restrictions to \(U\): compact inner regularity and compact-Hausdorff normality allow a continuous cutoff equal to one on any compact subset and supported in \(U\). Uniform approximation or the positive-functional description on these cutoffs then identifies the restricted measures. Hence \(\mu(u,\ldots)|U=\mu(w,\ldots)|U\). Replace slots successively to prove full mixed locality whenever corresponding tuples agree on \(U\). Every weak limit used a fixed continuous test with compact support, known domain membership and a common finite bound. No measure-order inference, increasing \(D\) closure, or plurifine claim occurs.

\section{n6. Exact one-factor gauge covariance, full weighted normalization/conductor rules, actual zero-arity curve degree and all reduction/degeneracy conventions.}
\textit{Supports: G331, n1, n2, n3, n4.}

For another proper primary \(\mathfrak b\), put \(d(v)=v_{\mathfrak a}(\mathfrak b)>0\), \(T(v)=v/d(v)\). G331 supplies a homeomorphism of sections and bounded continuous \(d,d^{-1}\). Transport degree-one homogeneous inputs by \(u_{\mathfrak b}(Tv)=u_{\mathfrak a}(v)/d(v)\). Rational profiles have the same ideals and denominators in both gauges, while tests satisfy \(v_{\mathfrak b}(J)(Tv)=v_{\mathfrak a}(J)(v)/d(v)\). The proposed measures have identical ideal moments, since the factors cancel:
\[
\mu_{\mathfrak b}(\boldsymbol u_{\mathfrak b})
=T_*\big[d\,\mu_{\mathfrak a}(\boldsymbol u_{\mathfrak a})\big]. \tag{GAUGE}
\]
Density proves this first on rational tuples. Uniform \(C\) and actual bounded decreasing \(D\) continuity extend it: transport preserves these convergences, and multiplication by fixed \(d\) and compact pushforward are weakly continuous. The factor is ONE \(d\), independently of dimension. The new mass is the old moment against \(v_{\mathfrak a}(\mathfrak b)\); for rational tuples it is \(e(I_1,\ldots,I_p,\mathfrak b)/\prod q_i\). The cocycle follows because successive weights multiply to \(v_{\mathfrak a}(\mathfrak c)\). On geometric atoms, \(\operatorname{ord}_Z(\mathfrak b)=d(\nu_Z)b_Z\) gives the same law directly.

Let \(S\) be the finite semilocal normalization, \(q\) its normal local centers above \(\mathfrak m\), \(f_q=[\kappa(q):\kappa(\mathfrak m)]\), and \(\rho_q\) restriction on the full boundary normalized by \(\mathfrak aS_q\). A regular conductor element kills \(S/R\). Any prime in its support has dimension below \(n\), since a chain above it extends strictly down to a minimal prime not containing that element. N1's exact-sequence top-form additivity therefore gives the same degree-\(n\) length form for \(R\) and \(S\). For each primary product \(P\), Artinian decomposition and composition series give
\[
\ell_R(S/PS)=\sum_q f_q\,\ell_{S_q}(S_q/PS_q).
\]
The residue degrees are field dimensions, including inseparable extensions. Comparing top forms yields
\[
e_R(I_1,\ldots,I_n)
=\sum_{\dim S_q=n}f_q\,e_{S_q}(I_1S_q,\ldots,I_nS_q).
\]
Testing on the dense ideal lattice, then extending by the two continuities, proves
\[
\mu_{R,\mathfrak a}(\boldsymbol u)
=\sum_{\dim S_q=n}f_q\,\rho_{q,*}
 \mu_{S_q,\mathfrak aS_q}(\rho_q^*\boldsymbol u). \tag{NORMAL}
\]
At lower-dimensional centers the degree-\(n\) contribution is ZERO; it is not their own lower-arity intrinsic measure. They and every support stay in the spaces and potential domain. There is no point deletion, nor an invented generic degree for a vertical component.

G331 describes the FULL normalization quotient: conductor-finite fibers are singletons; conductor-infinite fibers are exactly equal restrictions through \(S/\mathfrak c\to R/\mathfrak c\). It supplies positive \(C,D\) descent precisely for full-fiber-constant functions. Thus any compatible upstairs tuple in these classes gives actual downstairs inputs, and (NORMAL) computes their product by pushing along the full quotient, adding any masses on identified fibers. Both strata, nonminimal supports and mixed-zero patterns remain. This is a normalization/conductor rule, not arbitrary finite-map or envelope functoriality.

For an excellent reduced curve \(n=1,p=0\), normal centers are DVRs. A semivaluation there is a positive multiple of its order (nonzero elements are a uniformizer power times a unit); its only allowable support is zero. A regular conductor element has finite value on every branch, so the entire downstairs boundary is the finite branch set
\[
\nu_q=\operatorname{ord}_q/b_q,\qquad b_q=\operatorname{ord}_q(\mathfrak a)>0.
\]
The independently geometric zero-input measure is
\[
\mu_{R,\mathfrak a}()=\sum_q f_qb_q\,\delta_{\nu_q},\qquad
\mu_{R,\mathfrak a}(X)=e_R(\mathfrak a). \tag{CURVE}
\]
Here branches mean normal LOCAL centers, including multiple branches of one irreducible component.

Its actual Cartier comparison can also be proved directly on every curve primary model. Since \(IS\) is invertible, \(\operatorname{Spec}S\to Y\) exists. It is finite by its closed graph in \(\operatorname{Spec}S\times_RY\), and surjective because its finite image contains every component generic point of reduced \(Y\). It is the finite normalization. For effective exceptional Cartier \(E=(g)\) at a closed point, put \(B=\mathcal O_{Y,y}\), \(\widetilde B\) its finite semilocal normalization, \(C=\widetilde B/B\) of finite length. The snake sequence for multiplication by regular \(g\) gives
\[
0\to C[g]\to B/gB\to\widetilde B/g\widetilde B\to C/gC\to0.
\]
On finite-length \(C\), the kernel and quotient have equal lengths. Therefore the local Cartier degree, retaining residue weights, equals the sum of DVR orders with their residue degrees. Summing over \(y\), extending by effective differences and real coefficients, gives
\[
Q_Y(F)=\sum_q f_q\operatorname{ord}_q(F)
=\int_X H_F\,d\mu_{R,\mathfrak a}().
\]
This is an actual supported Cartier degree, agrees with n3--n4, and is independent of models. DVR lengths give \(\int v(J)\,d\mu=e_R(J)\). Zero-arity locality and continuity are constant identities.

Residue weights cannot be replaced by generic degree. Let \(k\subset l\) be ANY finite extension of degree \(d>1\), possibly inseparable, and \(R=k+t\,l[[t]]\subset S=l[[t]]\). If \(1,\alpha_2,\ldots,\alpha_d\) is a \(k\)-basis, \(R=k[[t]]\oplus\bigoplus_{i=2}^d t\,k[[t]]\alpha_i\) is a finite algebra over the excellent DVR \(k[[t]]\), hence excellent Noetherian. It is local, with \(\mathfrak m=tS\), residue \(k\), fraction field \(l((t))\), finite normalization \(S\), conductor \(tS\), generic degree one and residue degree \(d\). For \(J=t^rS\), a residue basis gives
\[
\ell_R(R/J^N)=1+d(rN-1),\qquad e_R(J)=dr.
\]
For gauge \(\mathfrak m\), (CURVE) is \(d\delta_{\operatorname{ord}_t}\), with the correct moment \(dr\), regardless of separability.

If \(n=0\), the reduced local ring is a field, its centered boundary is empty, and its only measure is zero; no negative-arity operator or degree-zero multiplicity mass is asserted. For \(n\ge2\), a zero input (represented by a unit) annihilates the product; a unit test has zero moment. For \(p=0\) there is no input zero to annihilate the reference measure. The tuple \(-t_i v(\mathfrak a)\) has mass \((\prod_i t_i)e_R(\mathfrak a)\), with empty product one. A zero ideal has value \(+\infty\), gives an unbounded bottom profile, and is outside the bounded product domain and outside primary gauges; it must not be multiplied formally by scalar zero. Define a zero scalar input as the actual zero function. Unit or zero ideals have no level-one section.

The reduction convention is necessary. \(k[[t]][\varepsilon]/(\varepsilon^2)\) and \(k[[t]]\) have the same semivaluation boundary, but parameter-power lengths are \(2s\) and \(s\). Nonreduced generic lengths require extra data. We claim reduced multiplicities only, while the genuine nilpotent thicknesses of Cartier divisors on reduced models were fully retained in n3.

\section{n7. Exact BFJ smooth OPERATOR specialization only in its published closed affine algebraically closed characteristic-zero range, including net topology.}
\textit{Supports: G331, n1, n2, n6.}

We identify the OPERATOR at precisely BFJ's published setting: a smooth CLOSED point of an affine variety over an algebraically closed characteristic-zero field, dimension \(n\ge2\), initially with gauge \(\mathfrak m\). G331 grants the full cone specialization on all semivaluations, with bottom adjoined; bottom is not bounded. No smooth multiplier statement is applied to singular or arbitrary-characteristic rings.

The complete BFJ Section 4 was read, not only the sequence statement 4.8. Its Cartier pairing is symmetric multilinear and has the ordinary exceptional intersection on a determination. Equations (4.3), Propositions 4.6--4.7 give ideal moments using the generic parameter COLENGTH convention. This equals our polynomial mixed multiplicity in the regular local ring, as follows.

A superficial parameter tuple exists by consecutive 17.2.2 choices and 17.3.1's joint-reduction construction. Joint reductions form a nonempty Zariski open set of coefficient tuples: for a fixed sufficiently large diagonal multi-Rees degree, the joint-reduction equation is surjectivity of a finite-dimensional multiplication matrix whose entries are linear in these coefficients, an open rank condition. Nakayama lifts residue surjectivity to the ideal equation. Conversely any joint reduction satisfies this equation in some such degree. Complete 17.4.9 says every joint reduction's parameter ideal has multiplicity equal to the polynomial mixed coefficient, at arbitrary Noetherian hypotheses. We use this ordinary published application here after reading its full statement/proof; it is not needed for the general monotonicity proof. In a regular local ring the parameter tuple is a regular sequence, so its Samuel multiplicity is exactly its quotient length. Thus generic parameter colength equals our mixed coefficient. Equivalently, for a superficial tuple repeated 17.4.6 has no annihilator correction by the CM property. These arguments are confined to the smooth range.

Hence BFJ 4.7 gives exactly n2's moments for rational primary profiles, with the same \(p=n-1\) arity and sign. Scaling, unit/zero conventions and lattice density give equality as Radon measures. For arbitrary \(C\) tuples, BFJ 4.4's smooth addition/order laws give the same negative-constant shift estimate (U). The mass is bounded by \(M^pe(\mathfrak m)\); uniform rational approximation consequently identifies its moments and measure with ours.

For bounded \(D\) tuples choose any actual decreasing \(C\)-representations. At each divisorial valuation their associated nef Weil coefficients decrease to the coefficients of the G331-identified BFJ limit, so the product net converges in the coefficient topology. BFJ 4.4's FULL proof gives upper semicontinuity in this topology: the strict coefficient inequality below a fixed nef Cartier divisor is an open finite-determination condition; Cartier majorants nearly attaining the infimum can be scaled slightly toward zero to make the inequalities strict. This provides open neighborhoods with intersection below any prescribed larger value. Together with intersection monotonicity, it proves continuity along ANY decreasing coefficient net: the values stay above the limit by order, while upper semicontinuity makes their limsup at most the limit. It needs no cofinal sequence. Apply this with the fixed ideal Cartier test \(Z(J)\). Positive ideal moments change the sign, and they converge exactly to our scalar limits in n2. Mass bounds and lattice density identify the entire measures on all semivaluations. The source's 4.8 sequence assertion alone was not promoted to a net theorem.

This is exact published smooth operator specialization, with the actual proof-topology argument. Remark 4.5's essential multiplier approximation for additivity remains a disclosed smooth dependency, as do resolution, Rees/Ramanujam intersections and usual CM multiplicity foundations. It supplies no general singular theorem. For other primary gauges use our proved (GAUGE); we do not attribute an unstated gauge theorem to BFJ. Smooth DVR and empty dimension zero follow directly from n6.

\section{n8. Reproducibility and diagnostic limits: all 54 bound files match, five source pages inspected, 6765 exact Newton checks pass; no computation certifies the proof.}
\textit{Supports: G331, n1, n2, n3, n4, n5, n6, n7.}

The reconstructed route retains the following diagnostics. Length-polynomial domination alone does not order mixed coefficients, as \((s-t)^2\) shows. Uniform convergence of Cartier evaluation functions alone is insufficient for intersection convergence; n4 supplies the necessary fixed-model identity (S) and polynomial convergence. The full normalization conductor is nonprimary and can have infinite value, so its quotient by an approximation exponent is not a uniform error; the fixed PRIMARY \(\mathfrak c_0\) is essential. Positive ideal-moment order is not measure order. Inferring locality from Riesz positivity or weak limits alone is invalid; n5 supplies the actual ratio-subideal argument and two compact-cover limits. Arbitrary increasing \(D\) closure and flattening nested \(D\)-representations were never assumed.

There was a presentation failure in the first owned file write: ordinary JavaScript string escaping removed inline TeX backslashes. The affected n1/n2 source files were completely rewritten using literal quoted data before delivery. This was an artifact encoding failure, not a failed mathematical theorem or evidence of one. The final integrity audit checks control characters and proof-body inclusion.

The finite diagnostic is an observation, separate from the universal proofs. The owned standard-library program \texttt{exact\_diagnostics.py} reads \texttt{bounded-evidence-input.json}. Axis exponents are \(1,2,3,4\), with optional generator absent or \((1,1),(1,2),(2,1),(2,2)\), giving 41 distinct monomial ideals after removing dominated generators. It uses exact rational primitive Newton normals, edge lengths, polygon areas, and evaluated monomial gauges. It checks the arithmetic Newton-facet formulas for 41 diagonals and every ordered pair of these ideals: 1681 symmetry, 1681 mixed-area, 1681 product-atom, and 1681 strict-maximum checks. All 6765 checks returned PASS, with zero failures. These computations do not prove the universal geometric interpretation of these formulas or any arbitrary-ring theorem; n1--n7 prove the operator without relying on this diagnostic. Nine residue arithmetic cases use \(d=2,3,5\), \(r=1,2,4\), \(N=1,2,3\); their first differences equal \(dr\), consistent with the universally proved curve example.

The owned byte-only \texttt{audit\_inputs.py} verified all 54 supplied files (21 mathematical DATA and 33 primary files) against the explicit manifests, with every size and SHA-256 matching. It copied the four qualified G331 components unchanged into \texttt{attachments/}. Input integrity establishes no mathematical correctness. The complete used immediate Stacks units were read: 0BEL 33.45.1--13, 0AG7 30.14.3, 01OF 31.33.1--14, 0AYA 42.26.1--4, 02QL 42.7.1, and 0B5S 30.16.1--3. The newly used algebra and smooth units and their exact application boundaries are in \texttt{source-application-ledger.md}; historical dependencies remain ordinary published background, not recursively freshly proved foundations.

The owned \texttt{inspect\_primary.py} rendered five exact bound PDF pages. Their pixels were inspected: Swanson--Huneke PDF 362/367/369 (printed 348/353/355) for the radical/colon hypotheses, coefficient convention, and full reduction formula; BFJ PDF 32/33 (printed 480/481) for the ideal sign/colength convention and coefficient-topology proof. The text extraction was not treated as a printed error. Source images and exact commands are retained in \texttt{reports/source-images/} and \texttt{reports/primary-image-log.json}.

No supplied executable was run or imported. No agent, native model job, unbound campaign inspection, hot feedback, account/runtime mutation, editor mutation, catalogue operation or publication occurred. No supervised heavy-compute facility was exposed and none was needed; only bounded local arithmetic and page rendering ran. Compilation is reserved to ROOT and was not run. Environment details, inputs and outputs are retained locally. Rerun from this directory: \texttt{python audit\_inputs.py}; \texttt{python inspect\_primary.py}; \texttt{python exact\_diagnostics.py}; \texttt{python build\_delivery.py}; \texttt{python final\_audit.py}. Execution, rendering and integrity checks do not establish mathematical correctness.

\section{n9. Complete ORIGINAL P12.4 solution candidate relative to entire qualified G331, with all proofs, exact scope and unchanged grant attachments; independent ROOT review remains.}
\textit{Supports: G331, n1, n2, n3, n4, n5, n6, n7, n8.}

We claim a complete solution candidate for the exact ORIGINAL P12.4, using only the entire unchanged qualified G331 as incoming support and the disclosed ordinary published dependencies. This is an unverified mathematical draft, for subsequent independent whole-original judgment by ROOT, not acceptance or publication.

The theorem is as follows. For an excellent local germ take its reduction \(R\), retain the entire intrinsic centered semivaluation boundary, and choose any proper primary section \(X_{\mathfrak a}\). If \(n=\dim R\ge1\), there is a unique positive symmetric real-cone multilinear mixed operator of arity \(p=n-1\) on the bounded independently geometric G331 class \(D\). On a real negative-effective curve-nef primary-model Cartier core it is the independently defined geometric atomic measure (A), with exact supported Cartier pairing (G), invariant under common primary refinements, normalization and all presentations giving the same functions. On rational ideal profiles its moments are \(e_R(I_1,\ldots,I_p,J)/\prod q_i\), and its reference mass is \(e_R(\mathfrak a)\) for the tuple of \(p\) copies of \(-1=-v(\mathfrak a)\). For general bounded inputs its mass and moments obey n2's bounds and extensions.

The operator is weakly continuous under uniform convergence inside the bounded domain and under ACTUAL bounded decreasing nets with prescribed limit in \(D\). It is local on EVERY open set where the corresponding bounded tuples agree. Primary changes obey the exact one-factor law (GAUGE), with the homogeneous input cocycle and the new test-reference mass. Finite normalization obeys (NORMAL) with all closed residue-field degrees, including inseparability, and zero top-degree contribution from lower-dimensional centers. Full conductor quotient fibers and every support remain. In dimension one the actual excellent reduced-curve Cartier reference measure is (CURVE); in dimension zero the boundary and measure are empty. Unit, zero, bottom and nilpotent conventions are as proved in n6. The operator agrees with the entire published BFJ Section 4 operator at its smooth affine CLOSED algebraically closed characteristic-zero hypotheses, and direct smooth DVR/empty conventions.

The proof is complete in n1--n7: n1 reconstructs the whole algebraic kernel, including unused labels, closure invariance, same-quotient all-field monotonicity, stabilized colons and top-form additivity; n2 proves positive function-dependent Riesz measures, every bounded representation comparison, and actual decreasing-NET continuity; n3 independently constructs the supported Artinian-thickening pairing, actual Koszul symmetry, divisor-CYCLE projection and raw ideal equality; n4 independently constructs positive atomic geometry and proves refinement invariance BEFORE the fixed-primary section-ideal bridge; n5 proves strict-maximum and full bounded open locality by legitimate compactly supported tests and two compact-cover limits; n6 proves every gauge/normalization/conductor/curve/degeneracy rule; n7 proves exact smooth OPERATOR specialization and net topology. N8 separates reproducibility observations and rejected shortcuts from proofs. No remaining P12.4 mathematical gap is claimed.

The theorem is relative to G331's directed PRIMARY-model geometric cone. It does not assert cofinality among all proper birational models, resolution in positive characteristic, a general nilpotent intersection theory, an envelope theorem, P12.5 energy, or P12.6 full arbitrary-finite-map calculus. Literal classical/CLD logarithmic radial extension comparison stays separate; the mixed-zero nodal obstruction is preserved. These are exact scope boundaries, not omitted P12.4 obligations.

The standalone \texttt{P12.4-paper.tex} contains the full new proof bodies, clean theorem and statement of the unchanged original target separately from its augmented reconstruction assignment. Its literal appendices include the whole bound G331 paper, annotation, source ledger and exact scope, supplied unchanged as data in \texttt{attachments/}. Other deliverables are \texttt{source-application-ledger.md}, \texttt{exact-scope.md}, \texttt{whole-proof-coverage.md}, \texttt{original-scope-and-noncircularity-decision.md}, \texttt{native-notes.json}, all \texttt{n1.tex} through \texttt{n9.tex}, and owned bounded inputs/scripts/reports. The paper does not hide kernel proofs behind incoming working sketches. Rerun commands and evidence limits are in n8. Compilation, final acceptance and publication remain with ROOT. Candidate=true is solely this full original-solution claim, subject to later independent review.

\appendix
\section{Entire qualified G331 source, unchanged literal DATA}
The following complete bound source is supplied unchanged as an attachment.
Its no-operator qualification is preserved; the NEW proofs above supply the
intersection and locality theorems rather than importing them.
\begin{Verbatim}[breaklines,fontsize=\tiny]
\documentclass[11pt]{article}
\usepackage[T1]{fontenc}
\usepackage{lmodern}
\usepackage{amsmath,amssymb}
\usepackage[margin=1in]{geometry}
\usepackage[hidelinks]{hyperref}
\usepackage{microtype}
\setlength{\parindent}{0pt}
\setlength{\parskip}{0.6em}
\emergencystretch=4em
\title{A geometric valuative psh cone on excellent local germs\\
Actual Cartier positivity, uniform comparison, and full conductor descent}
\author{Fresh local reconstruction of ORIGINAL P12.2}
\date{8 October 2026}
\begin{document}
\maketitle
\begin{abstract}
On the intrinsic centred semivaluation boundary of an excellent local germ,
we define a geometric core by actual negative effective real Cartier divisors
on actual primary blowups, with nonnegative degrees on all contracted integral
curves. We prove its continuous uniform closure and actual decreasing completion
equal the corresponding ideal classes. A fixed primary correction to finite
global sections, corrected Keeler generation on the full closed fibre, and
explicit real ample perturbation prove the comparison without a normality,
smoothness or characteristic restriction. We reconstruct full finite-injection
and normalization/conductor descent, including mixed-zero and bottom patterns,
and recover the full smooth BFJ class on all semivaluations at its exact published
hypotheses. A nonisolated cusp separates the entire finite-real raw cone from its
completion. This is an unverified complete-solution candidate for P12.2;
independent whole-original review remains required.
\end{abstract}
\tableofcontents
\newpage
\section*{Main theorem and proof map}
Let $(R,\mathfrak m)$ be excellent local and replace it by its reduction.
For every proper primary $\mathfrak a$, let $V_{\mathfrak a}$ be the compact
section of centred real semivaluations with $v(\mathfrak a)=1$, retaining every
prime support. Define $Q$ by finite nonnegative real sums of negative
unit-or-primary ideal values, $C$ by uniform closure, and $D$ by actual infima
of nonempty decreasing $C$-nets, with bottom allowed.

On actual primary blowups, take actual negative effective finite real Cartier
presentations having nonnegative degree on every contracted proper integral
curve. Their evaluations define $M_{\rm geom}$ independently of ideal
positivity. Then
\[
 Q\subset M_{\rm geom}\subset C,\qquad
 C_{\rm geom}=C,\qquad D_{\rm geom}=D.
\]
For every reduced finite injection, the pullback image of either completed
class is precisely its full-fibre-constant subclass. Normalization fibres are
singletons at finite conductor value and precisely conductor-quotient fibres
at infinite value. Primary changes obey the exact homogeneous cocycle.
At a smooth affine closed germ over an algebraically closed characteristic-zero
field, $C$ is exactly the finite continuous BFJ subclass and $D$ is the full
BFJ all-semivaluation extension with the identically bottom function adjoined.
Smooth DVR and empty dimension zero are direct.

The proofs are the numbered local chain below: n1 gives the boundary; n2--n6
give full algebraic descent; n7 discharges the exact published geometric inputs;
n8 constructs and compares the independent cone; n9 gives full BFJ recovery;
n10--n11 preserve raw-cone and radial distinctions; n12 records the exact
original-target decision. Each support list contains only earlier local IDs.
Source applications are ordinary published applications, with disclosed
recursive historical dependencies, and no supplied draft is an accepted premise.
Separate envelopes, products, energy and their geometric operator theory are
the P12.3--P12.6 targets. They are not claimed here.
\section{n1. Boundary, finite lifting, conductor fibres, and completion}
\textit{Earlier local supports: none.}

All rings below are commutative. For a Noetherian local ring \((R,\mathfrak m)\), a centred real semivaluation is a map to \([0,+\infty]\) satisfying \(v(0)=+\infty\), \(v(1)=0\), multiplicativity, the ultrametric inequality, and \(0<v(\mathfrak m)<+\infty\). Its support \(v^{-1}(+\infty)\) is a prime strictly below \(\mathfrak m\); no minimal-support restriction is imposed. Ideal values are minima over finite generating sets. Relations between generating sets prove independence and
\[
v(IJ)=v(I)+v(J),\qquad v(I+J)=\min\{v(I),v(J)\}.
\]
Units have value zero. A proper primary ideal always means primary for the marked maximal ideal.

Fix a proper primary ideal \(\mathfrak a\). The section \(V_{\mathfrak a}(R)=\{v:v(\mathfrak a)=1\}\) is compact Hausdorff. Indeed semivaluation axioms and \(v(\mathfrak m)=1\) are closed conditions in the compact product \([0,+\infty]^R\); addition, minimum and the order relation have the required continuity/closedness. Choose \(\mathfrak m^r\subset\mathfrak a\subset\mathfrak m^s\), with \(r\ge s\ge1\). Then \(s v(\mathfrak m)\le v(\mathfrak a)\le r v(\mathfrak m)\). Rescaling identifies the two sections. Every primary ideal value is finite positive continuous and bounded there. For another primary ideal \(\mathfrak b\),
\[
H_{\mathfrak a,\mathfrak b}(v)=v/v(\mathfrak b),\qquad
u_{\mathfrak b}(H_{\mathfrak a,\mathfrak b}v)=u_{\mathfrak a}(v)/v(\mathfrak b).
\]
The denominator and its inverse are bounded. Substitution proves the inverse and the exact three-gauge cocycle. The centred cone is the product of its section and the positive radial coordinate \(v(\mathfrak a)\); hence these sections represent its intrinsic ray quotient. Nonprimary proper ideals give only their open finite-value domain, which can be noncompact. Units and zero ideals have no level-one section.

Nilpotents have infinite value. Comparing \(f+n\) and \((f+n)-n\) proves that reduction preserves every semivaluation and ideal gauge. A zero-dimensional reduced local ring is a field and has empty centred boundary. For a positive-dimensional excellent reduced local ring, a penultimate prime of a saturated chain to the maximal ideal gives a one-dimensional excellent domain quotient. Its finite normalization has a local DVR factor, whose order pulls back and normalizes into the section. This proves nonemptiness. Ordinary background used throughout is lying over, prime avoidance, finite Artinian decompositions, and the quotient/finite-type-localization stability, Nagata finiteness and universal catenarity consequences of excellence. None is used to assert resolution in arbitrary characteristic.

Every finite integral injection \(R\hookrightarrow B\), with \(B\) reduced and nonzero, gives a continuous surjective restriction \(V_{\mathfrak aB}(B)\to V_{\mathfrak a}(R)\). Here a finite algebra has the finite union of its closed-centre sections above \(\mathfrak m\). To lift \(v\) with support \(\mathfrak p\), choose a prime \(P\) of \(B\) above \(\mathfrak p\) and form the finite field extension \(\operatorname{Frac}(B/P)/\operatorname{Frac}(R/\mathfrak p)\). Extend the valuation ring by taking the integral closure, a prime over its maximal ideal, and a dominating valuation ring. The domination assertion follows from Zorn: in a maximal dominating local subring, either adjoining \(x\) or \(x^{-1}\) leaves the maximal ideal proper. If \(1=\sum a_i x^{-i}\) with all \(a_i\) in that ideal, multiplying by a power and dividing by \(1-a_0\) makes \(x\) integral, so lying over makes the ideal proper after adjoining \(x\). Localizing contradicts maximality unless one of the two elements already belongs. Intersection with the old field is the original valuation ring: an element outside it has inverse in its maximal ideal. An algebraic equation for any new nonzero field element has two least-valued terms; thus an integer multiple of its value is in the old value group. Dividing this old real value by the integer gives a well-defined order-preserving real embedding, using common multiples. Integral elements have nonnegative value by the least-term test for a monic equation. Compose with \(B\to B/P\), assigning \(P\) infinite value. This construction first kills the support; it never extends a nonzero-support semivaluation to the original fraction field. Compactness makes restriction a quotient map.

Let \(S\) be the finite normalization of an excellent reduced \(R\), and \(\mathfrak c=(R:S)\). Clearing denominators of finitely many generators of \(S\) gives a regular conductor element. The conductor is an ideal of both rings and
\[
R=S\times_{S/\mathfrak c}(R/\mathfrak c).
\]
An element of \(S\) whose residue comes from \(R\) differs from an element of \(R\) by an element of \(\mathfrak c\), proving this identity. The normalization boundary is the finite disjoint union of local normal centres above \(\mathfrak m\). Idempotents choose a minimal component; localization at its centre changes no values of denominators. Chinese remainder elements separating the finitely many centres, and their complements, have open positive-value loci selecting these pieces. The pieces are therefore clopen. Branches mean these local centres, not merely minimal primes.

If \(v(\mathfrak c)<+\infty\), choose \(c\in\mathfrak c\) of finite value. Every lift satisfies \(w(s)=v(cs)-v(c)\), so it is unique. The subtracted number is finite. If \(v(\mathfrak c)=+\infty\), all lifts kill the conductor, and equality after restriction \(S/\mathfrak c\to R/\mathfrak c\) is exactly their identification. This infinity stratum is closed; its finite complement is open, without a closedness claim. A fibre-constant continuous function descends continuously by the compact quotient. A fibre-constant extended nonpositive upper semicontinuous function descends upper semicontinuously: its closed superlevel sets have compact, hence closed, images. These assertions alone supply no positivity descent.

Completion preserves the entire boundary for every Noetherian local ring. For \(F\in\widehat R\), choose representatives \(F_N\) modulo \(\mathfrak m^N\), and put \(d=v(\mathfrak m)>0\). Differences have value at least \(Nd\). If some \(v(F_N)<Nd\), later values stabilize; otherwise they tend to infinity. Define \(\widehat v(F)\) by this limit. It is independent of representatives and satisfies
\[
\min\{\widehat v(F),Nd\}=\min\{v(F_N),Nd\}.
\]
Applying this to sums and products proves all axioms and uniqueness, including infinite values; completed units have unit representatives. Compose with \(q(t)=t/(1+t)\), \(q(+\infty)=1\). The displayed continuous truncations have uniform error at most \((1+N/r)^{-1}\), because \(d\ge1/r\). Thus completed evaluation is continuous and completion restriction is a homeomorphism. Every completed primary \(J\) contains \(\widehat{\mathfrak m}^N\). Representatives modulo that power of its generators belong to \(J\cap R\) and, with the boundary power, generate \(J\). Hence \(J=(J\cap R)\widehat R\). The quotient isomorphism also proves contraction of an extended primary ideal recovers it. Unit-or-primary profiles and their uniform and decreasing classes therefore agree under completion. No equivalence of geometric model categories is asserted.

This is a mathematical argument with the stated ordinary background. Swanson--Huneke 6.4.1--2 were read in full and agree with the domination proof. Artifact n1.md is authoritative and is reproduced in argument.md and the paper. Rerun: python audit\_inputs.py; python build\_delivery.py. Execution does not verify the mathematics.

\section{n2. Exact algebraic applications and their qualifications}
\textit{Earlier local supports: n1.}

Use n1 for boundary conventions. The following are ordinary published theorem applications at their stated hypotheses, with their immediate units checked from the bound primary data. This is not an independent reconstruction of all historical proofs.

Swanson--Huneke is the current author PDF created July 2024, not asserted byte-identical to the 2006 printing. Its complete 6.8.1--4 statements and proofs give the following consequence. Let \(E\) be Noetherian semilocal, \(J\) contain a boundary power, and \(Q\) be a finite ideal. If every normalized centred semivaluation satisfies \(w(Q)\ge w(J)\), and \(Q\) has zero image on every zero-dimensional minimal component where \(J\) is nonunit, then
\[
Q\subset\overline J.
\]
Apply 6.8.4 on each positive-dimensional local domain component, normalizing each centred rank-one DVR by the boundary ideal. Apply 6.8.3 to all minimal components. On zero-dimensional components there are no boundary points; the extra zero-image condition is needed. In all uses below \(Q=K(\mathfrak aE)^L\), with \(L>0\), so this condition holds. Unit components of \(J\) impose no condition. For nonreduced \(E\), lift the componentwise monic equations, multiply them over the finitely many minimal primes, then raise to a power to kill the nilpotent error. The resulting equation is monic with degree-i coefficients in \(J^i\). This is the finite-minimal-prime argument of 1.1.5, also read completely. PDF inspection of printed page 140 confirmed the integral-closure bars omitted in the text extraction.

The common exponent needed later has a short direct proof. Each \(qt\), \(q\in Q\), is integral over \(E[Jt]\) with its homogeneous ideal-integrality equation. Thus \(E[(J+Q)t]\) is finite graded over \(E[Jt]\). Choose homogeneous module generators of degrees at most \(h\). Its degree \(N+h\) terms have coefficients in \(J^{N+h-d}\subset J^N\). Consequently
\[
Q^{N+h}\subset (J+Q)^{N+h}\subset J^N\quad(N\ge1).
\]
This works for nonreduced rings and assumes no finiteness of the full normalized Rees algebra. The complete 1.2.3--5 reduction chain was checked too; printed page 6 confirms the criterion uses integral closure.

For an excellent local domain \(A\) of positive dimension and a nonzero proper primary ideal, 10.2.2, including its proof ending at 10.2.3, constructs finitely many Rees DVRs centred at the closed point. The full 10.4.3 proof gives their residue transcendence degree \(\dim A-1\) when \(A\) is formally equidimensional. Excellence supplies universal catenarity, and the full B.5.1 statement/proof identifies this with local formal equidimensionality for Noetherian domains. B.3.2 and its immediate context give, for a finite-type birational domain \(B\) over excellent local domain \(A\), and \(P\cap A=\mathfrak m\),
\[
\dim B_P+\operatorname{trdeg}_{A/\mathfrak m}\kappa(P)=\dim A.
\]
The book describes B.3.2's proof as a modification of an earlier dimension argument. Its recursive historical dependencies, including Mori--Nagata and Ratliff, are disclosed rather than claimed independently read.

Stacks 0C23, full immediate HTML lemma/proof, states that normal Noetherian local rings with normal formal fibres have normal completion; excellence suffices. Its proof refers to Algebra 10.163.8, not recursively read. A normal local completed ring is a domain. Thus an excellent normal local domain is analytically irreducible and analytically unramified. This is not claimed for an arbitrary nonnormal domain.

Exact Rond--Spivakovsky arXiv 1404.4885v1, complete pages 1--2 definitions and statements 1.3, 1.4 and 1.7, yields an order bound at precisely this analytic domain. For a centred divisorial valuation \(\nu\) and the finitely many maximal-ideal Rees valuations \(\nu_j\), 1.3 gives \(\nu(f)\le c_j\nu_j(f)\). Set \(C_0=\max_j c_j\nu_j(\mathfrak m)\). Statements 1.4 and 1.7 imply
\[
\nu(f)\le C_0\min_j\frac{\nu_j(f)}{\nu_j(\mathfrak m)}
=C_0\overline{\operatorname{ord}}_{\mathfrak m}(f)
\le C_0(\operatorname{ord}_{\mathfrak m}(f)+K).
\]
For nonzero nonunits the order is at least one and finite by Krull intersection, giving \(\nu(f)\le C_0(1+K)\operatorname{ord}_{\mathfrak m}(f)\). Units have both values zero and are handled separately. The literal strict Remark 1.6 fails on units and is not used. Swanson's full Rees-valuations Section 5 was read; its 5.2 is stated for complete domains. The general-domain application here uses exact Rond--Spivakovsky 1.3. Original Rees 1989 and the historical Izumi proofs have not been independently read. Published applications with these disclosed dependencies are allowed by the task; metadata and old reading receipts are not grants.

Artifacts: n2.md, source-application-ledger.md, reports/source-images/book-140.png and book-6.png, reports/input-hashes.json. Rerun: python inspect\_primary.py; python audit\_inputs.py; python build\_delivery.py. No finite check certifies these published applications or the new consequences.

\section{n3. Exact cones, density, and normal-target traces}
\textit{Earlier local supports: n1.}

Use n1's sections. Define \(P_R\) to be the raw finite nonnegative REAL combinations of gauges \(v(I)\), where \(I\) contains a power of the primary boundary ideal, with the unit ideal allowed. Put \(Q_R=-P_R\), let \(C_R\) be its uniform closure in finite continuous functions, and let \(D_R\) consist of infima of nonempty decreasing \(C_R\)-nets. Include zero and the identically minus-infinity bottom. A finite semilocal algebra uses the product classes over its closed centres. These agree with global semilocal ideal profiles: local primary ideals contract to ideals primary at their own centre and unit at the others, and tuples glue in a common Artinian boundary quotient. Denominators synchronize by powers. Rational profiles \(-v(I)/k\) are uniformly dense in \(C_R\): approximate finitely many real coefficients rationally and turn their sum into a product-ideal profile. The class \(C_R\) is a closed convex cone, stable under maxima by ideal sums. Its decreasing completion is upper semicontinuous and stable under addition, positive scaling and finite maxima by product directed sets. We do not assume increasing-limit closure, arbitrary-infimum closure, or flattening of nested nets.

The rational difference span of bounded gauges is a lattice: after a common denominator every expression is \((v(I)-v(J))/k\), and
\[
\min\{v(I)-v(J),v(K)-v(H)\}=v(IH+KJ)-v(JH).
\]
It contains rational constants since \(v(\mathfrak a)=1\), and separates points using \(v((f)+\mathfrak a^N)=\min\{v(f),N\}\), and centre-separating ideals in the semilocal case. Its UNIFORM closure is a real vector lattice. The compact lattice approximation proof makes this closure all real continuous functions: affine combinations of a separating function match a target at two points; rational approximations preserve strict local error tests. For a fixed first point take finitely many minima to bound the target from above throughout the space, while matching near that point; then finitely many maxima over first-point neighbourhoods give a uniform approximation. Negating interchanges the needed bounds. No integer ideal formula is asserted for an arbitrary finite irrational span.

Let \(A\) be excellent normal local domain and \(A\subset B\) a finite domain extension, with fraction fields \(K\subset L\), \(n=[L:K]\). For a \(B\)-ideal \(J\), choose finite \(A\)-module generators \(b_i\) and define \(\mathcal N(J)\) as the coefficient ideal of
\[
N_{L(T)/K(T)}\Big(\sum_i T_i b_i\Big)
=\prod_{\sigma:L\hookrightarrow M}\Big(\sum_i T_i\sigma(b_i)\Big)^q,
\qquad q\,|\Sigma|=n.
\]
Here \(M/K\) is a finite field containing all embedding images. The generic inseparable length \(q\) follows by decomposing \(L\otimes_K\overline K\) into its Artinian factors: determinant of multiplication on a length-q factor is its residue scalar to the power q. The coefficients are integral over \(A\) and belong to \(K\), hence belong to \(A\). Linear substitution in both directions when changing generators proves independence of this coefficient ideal. If \(J\supset\mathfrak a^N B\), add base generators \(x\) of \(\mathfrak a^N\) redundantly: the pure variable coefficient is \(x^n\). The coefficient ideal therefore contains a boundary power and is unit-or-primary.

Let \(E\) be the finite integral closure of \(A\) in \(M\). Lift any base semivaluation to \(E\) by n1, including nonzero support, and compose with the embeddings to obtain \(w_\sigma\) on \(B\). Minimum coefficient values of polynomials multiply. To prove this, quotient by the support, enter the valued quotient fraction field, divide each nonzero polynomial by a coefficient of smallest finite value, and reduce to nonzero residue polynomials. Their product is nonzero. A polynomial vanishing modulo support has infinite coefficient minimum and causes no difficulty. Applying this to the displayed norm polynomial gives on EVERY support
\[
v(\mathcal N(J))=q\sum_{\sigma\in\Sigma}w_\sigma(J).
\]
Boundary-power ideals make these sums finite. Define \(\mu_v=q\sum_\sigma\delta_{w_\sigma}\), of mass \(n\). Its integrals on all ideal gauges are fixed by the coefficient identity; the density just proved and equal mass make it independent of the lift on every continuous function, hence as a finite measure. Its support is the full restriction fibre. Indeed for any fibre point on \(B\), extend it to \(E\) by n1 with the identity embedding of \(B\). That point occurs with positive weight. Choice independence forces it into the support of every such measure. Thus the whole fibre is finite of cardinality at most \(n\). This proof uses generic inseparable lengths and quotient-support Gauss values; it does not assume a defectless ramification/residue-degree formula.

The trace \(TH(v)=\int H\,d\mu_v\) is positive, of norm \(n\), and continuous-valued. Continuity holds first on gauges by the coefficient identity, then on all continuous functions by uniform density and the mass bound. It preserves \(C\), mapping a rational negative profile to \(-v(\mathcal N(J))/k\). A finite positive sum commutes with decreasing nonpositive extended limits: if one summand goes to bottom so does the sum; otherwise directedness approximates all finite limits simultaneously. Consequently a fibre-constant \(D_B\)-function descends via \(T/n\) to an actual decreasing \(C_A\)-net, as does a fibre-constant \(C_B\)-function.

For any reduced finite injection into normal local \(A\), choose the minimal source components dominating \(A\), at least one of which exists over the generic point, and normalize them. Sum their traces divided by the sum of their generic degrees. At every base semivaluation all traced values equal the descended value by constancy on the FULL original fibre. Vertical components receive no fabricated generic degree, but remain constrained by that full constancy. This proves exact \(C,D\) descent for normal targets. Pullback of ideals and limits supplies the converse. Every lift has the same base primary-gauge denominator, so these assertions obey n1's gauge cocycle.

Artifacts: n3.md and the paper. Rerun: python build\_delivery.py. This coefficient trace is proved here rather than imported from a supplied label. It is not a construction of envelopes, products, energy, or a degree trace on arbitrary nonnormal targets.

\section{n4. Bounded descent for every reduced excellent finite injection}
\textit{Earlier local supports: n1, n2, n3.}

Using n1--n3, we prove
\[
f^*C_R=\{u\in C_B:u\text{ is constant on every restriction fibre}\}
\]
for every excellent reduced local \(R\) and nonzero reduced finite injection \(R\hookrightarrow B\). Pullback is an isometry. The theorem concerns the continuous UNIFORM cone, not the raw cone or a decreasing completion.

Induct on target dimension. Dimension zero has empty sections. Normal targets are n3. In dimension d, first prove normalization descent, using the arbitrary finite-injection theorem below dimension d; then prove arbitrary injections by normalization base change. Let \(S\) normalize \(R\), \(\mathfrak c\) be its conductor, \(C_0=R/\mathfrak c\), and \(E=S/\mathfrak c\). A regular conductor element gives \(\dim C_0<d\). The injection \(C_0\hookrightarrow E\) stays injective after reduction because nilpotence is detected in an injected subring. Restriction of primary gauges and uniform limits to the quotient section preserves the class. Thus for fibre-constant \(u\in C_S\), its continuous descent \(h\) from n1 has quotient restriction \(h_0\in C_{C_0}\) by induction. If the quotient section is empty, the conductor value is finite at every point and quotient estimates below are unnecessary.

Fix \(\varepsilon>0\). Choose rational profiles \(j=-w(J)/k\) and \(\kappa=-v(K)/k\) approximating \(u\) and \(h_0\) with error less than \(\varepsilon\). All ideals contain boundary powers. We can make \(J\) proper at every positive-dimensional centre: add an arbitrarily small negative boundary constant before rational approximation, or multiply a rational profile ideal by the boundary ideal with an arbitrarily large common denominator. Start with smaller errors to accommodate this correction. Synchronize denominators by ideal powers and enlarge k until \(1/k<\varepsilon\). On the entire quotient boundary,
\[
|w(K_E)-w(J_E)|<2k\varepsilon.
\]
Choose integer \(L\) with \(2k\varepsilon<L\le2k\varepsilon+1\). The ideal \(Q=K_E(\mathfrak aE)^L\) has \(w(Q)\ge w(J_E)\) everywhere. Its POSITIVE boundary factor kills its image on zero-dimensional minimal components. The exact n2 criterion therefore gives \(Q\subset\overline{J_E}\), on all minimal components and even a nonreduced quotient. The finite graded argument gives one \(h'\ge0\) with \(Q^{N+h'}\subset J_E^N\) for every N.

Put \(I_N=J^N\cap R\). It contains a boundary power and is proper, hence primary. The conductor fibre product in n1 gives the exact residue identity
\[
(I_N+\mathfrak c)/\mathfrak c=J_E^N\cap C_0.
\]
For the nontrivial inclusion, lift a residue in the right side to an element of \(J^N\). Since its residue comes from \(C_0\), that very element lies in \(R\), and hence in \(I_N\). Thus the image contains \(K^{N+h'}(\mathfrak aC_0)^{L(N+h')}\). At a conductor-infinite point its quotient semivaluation yields
\[
\frac{v(I_N)}{kN}\le(1+h'/N)\left(\frac{v(K)}k+\frac Lk\right).
\]
All terms here are finite. At a conductor-finite point the unique lift \(w\) and \(\mathfrak cJ^N\subset I_N\) yield instead
\[
\frac{w(J)}k\le\frac{v(I_N)}{kN}\le\frac{w(J)}k+\frac{v(\mathfrak c)}{kN}.
\]
The latter estimate is used ONLY at finite conductor value. On every point the lower bound from \(I_N\subset J^N\) gives \(v(I_N)/(kN)\ge-h-\varepsilon\). The limit upper bound is at most \(-h+4\varepsilon\): at infinity use \(-\kappa\le-h+\varepsilon\) and \(L/k<3\varepsilon\), and at finite conductor use \(-j\le-h+\varepsilon\).

The continuous functions \(b_n=v(I_{2^n})/(k2^n)\) decrease, since \(I_N^2\subset I_{2N}\). Their increasing strict-test open sets \(\{b_n+h<5\varepsilon\}\) cover the compact section. One n covers it entirely. Together with the common lower bound this proves \(\|-b_n-h\|_\infty<5\varepsilon\), and then \(h\in C_R\). These negative dyadic profiles INCREASE internally. They are not used as a decreasing regularization, and no increasing-limit closure is invoked.

For arbitrary \(R\hookrightarrow B\), form \(B'=(B\otimes_R S)_{\rm red}\). Base change of an integral surjective spectrum map is surjective, so its kernel in reduced \(S\) lies in every prime and is zero. Hence \(S\hookrightarrow B'\) is a finite injection. Pullback of \(u\) is fibre constant over each normal local centre of \(S\); equal restrictions there imply equal underlying restrictions to \(R\). N3 descends it to \(C_S\), exactly the pullback of \(h\). Normalization descent just proved gives \(h\in C_R\), closing induction. Ideal extension gives pullback inclusion and surjectivity in n1 gives the isometry.

Define \(N_R\) as the continuous functions whose normalization pullbacks belong to the RAW local negative real ideal cones and are fibre compatible. We have only
\[
Q_R\subset N_R\subset C_R,\qquad
\overline{Q_R}^{\|\cdot\|_\infty}=\overline{N_R}^{\|\cdot\|_\infty}=C_R.
\]
No raw equality is inferred. All conductor-infinite quotient supports remain. Artifacts: n4.md and the paper. Rerun: python build\_delivery.py. The incompatible-ideal contraction shortcut is deliberately absent; n11 diagnoses it.

\section{n5. Unbounded strictly negative descent by insertion}
\textit{Earlier local supports: n1, n2, n3, n4.}

Use n1--n4. A fibre-constant \(u\in D_B\) descends into \(D_R\) whenever its descent \(h\) is everywhere strictly negative, including partial minus-infinity and discontinuous unbounded functions. The descent is upper semicontinuous by n1, so on a nonempty compact section \(h\le-c\) for some finite \(c>0\). All-bottom is directly the limit of \(-N v(\mathfrak a)\); zero is separate. Normal targets follow from n3. Induct on dimension exactly as in n4: prove normalization from lower-dimensional finite-injection descent and then arbitrary injections by reduced normalization base change.

In the normalization step use \(S,\mathfrak c,C_0,E\) from n4. A decreasing representation of \(u\) restricts to the quotient, whose fibre-constant limit descends by induction to \(h_0\in D_{C_0}\). We prove an insertion lemma: for every finite continuous test \(F>h\) there is \(q\in C_R\) with \(h<q<F\). Choose a FINITE positive lower bound \(e\) for the lower semicontinuous gap \(F-h\), for example any number below half the minimum of 1 and its infimum. When that infimum is infinite, e remains finite; no infinity subtraction is performed.

From a decreasing continuous representation of \(u\), choose one term \(g\) satisfying \(u\le g<-c/2\) and \(g<F\circ\rho-e/2\). Each point eventually satisfies these strict continuous tests; their open sets and directedness give one common index by compactness. No uniform convergence to \(u\) is used. Choose \(0<b<1/2\) with \(b\|g\|<e/8\). The scaled term \((1-b)g\) is still below \(F-e/4\), and exceeds finite \(u\) by at least \(bc/2\). Uniform rational approximation with error less than \(\min(e/8,bc/4)\), further tightened to retain strict negativity, gives \(j=-w(J)/k\) and a fixed \(s>0\) such that
\[
u+s\le j<F\circ\rho-e/8.
\]
The first inequality is automatic at bottom. \(J\) is proper at all centres and contains a boundary power. Put \(M=\|j\|\).

We need a rational quotient profile \(\kappa=-v(K)/k\) with \(h_0\le\kappa\) and \(\kappa\circ\rho_E<j|_E+\delta\), for arbitrarily small \(\delta>0\). Here is the needed construction. Choose \(0<b_1<1/2\) with \(b_1M<s/4\), and a decreasing \(C_{C_0}\)-net \(g_i\) to \(h_0\). Since \(u_E\le j_E-s\), we have \((1-b_1)u_E<j_E-s/4\). Compactness on E selects one i with its scaled pullback below \(j_E-s/8\); simultaneously compactness on \(C_0\) selects \(g_i<-c/2\). Scaling therefore lifts it above \(h_0\) by a positive uniform gap at finite points. Rational approximation with error less than \(\min(b_1c/4,s/16)\) gives \(\kappa\) with the stated two bounds, even with a negative gap to \(j_E\). Synchronize denominators with J and enlarge k so \(1/k<\delta\). If the quotient boundary is empty omit this construction.

Choose integer \(L\) with \(k\delta<L\le k\delta+1\). On E the ideal \(Q=K_E(\mathfrak aE)^L\) dominates \(J_E\) in value. N2, including the zero-dimensional boundary factor, gives \(Q\subset\overline{J_E}\) and a common \(h'\) with \(Q^{N+h'}\subset J_E^N\). Contract \(I_N=J^N\cap R\). The exact residue identity of n4 gives at infinity-conductor points
\[
\liminf_N[-v(I_N)/(kN)]\ge\kappa-2\delta\ge h-2\delta.
\]
At finite-conductor points its finite-value estimate gives limit j, at least \(h+s\). The dyadic \(f_n=-v(I_{2^n})/(k2^n)\) increase, are finite continuous, and satisfy \(f_n\le j\) on every lift. Their strict open sets \(\{f_n>h-3\delta\}\) cover the compact section, including bottom points: these sets are open because h is upper semicontinuous. Select one n valid everywhere.

Choose \(\delta\) so small that \(\eta=3\delta/c<1/2\) and \(\eta M<e/16\), and put \(q=(1-\eta)f_n\in C_R\). At finite \(h\le-c\),
\[
q>(1-\eta)(h-3\delta)\ge h+\eta c-3\delta(1-\eta)=h+3\delta\eta>h.
\]
At bottom q is finite. Also \(q\le(1-\eta)j\le j+\eta M<F\). This proves insertion using real positive scaling, not adding a positive constant to the class.

Let \(\mathcal Q\) be all continuous \(C_R\)-functions strictly above h. Zero belongs. For two members apply insertion to their continuous minimum; the test need not itself be psh. This gives a member below both, so \(\mathcal Q\) is downward directed. At any v and finite \(t>h(v)\), upper semicontinuity and compact Hausdorff normality give a finite continuous F strictly above h with \(F(v)<t\): choose \(h(v)<\ell<t\), a neighbourhood where \(h<\ell\), and a Urysohn cutoff supported there, interpolating from \(\ell\) to a global strict upper bound outside. Insertion supplies \(q(v)<t\). Thus \(\inf\mathcal Q=h\), and reverse pointwise order is an ACTUAL decreasing \(C_R\)-net to h. This closes normalization and the finite-injection induction. Neither an increasing completion step nor a flattening theorem is needed.

Artifacts: n5.md and the paper. Rerun: python build\_delivery.py. This proof reconstructs strict-negative descent without accepting any supplied provenance label. Mixed-zero descent is next.

\section{n6. Divisorial density, normal zero dichotomy, and full completed descent}
\textit{Earlier local supports: n1, n2, n3, n4, n5.}

Use n1--n5. We prove the full statement
\[
f^*D_R=\{u\in D_B:u\text{ is constant on every restriction fibre}\}
\]
for every nonzero reduced finite injection over an excellent reduced local ring. It includes allzero, allbottom, mixedzero, partialbottom, conductor-infinite, and nonminimal supports.

First prove divisorial density on an excellent local domain \(A\). A basic neighbourhood of v in its primary section imposes finitely many strict coordinate tests. Choose a generator \(a_0\) of \(\mathfrak a\) with \(v(a_0)=1\) and put all \(a_j/a_0\) in a finite-type birational subalgebra \(B_0\subset\operatorname{Frac}(A)\). For each finite positive tested coordinate f choose rational bounds \(p/q<v(f)<r/s\) within its required interval, and put \(f^q/a_0^p\) and \(a_0^r/f^s\) in \(B_0\). For an infinity coordinate use only a lower bound and its first ratio. Units and zero have fixed values. Every denominator has finite value, so even a nonzero-support v extends nonnegatively to this algebra. Its positive-value prime P contains \(a_0\) and every strict-test ratio and contracts to \(\mathfrak m\). Thus \((B_0)_P\) is excellent domain of positive dimension. A Rees DVR of its maximal ideal, furnished by n2, is positive on P and nonnegative on \(B_0\). Its valuation \(\nu\) has \(\nu(\mathfrak a)=\nu(a_0)>0\); dividing by this number satisfies all the original tests.

This DVR is divisorial over A itself. N2's birational dimension formula gives \(\dim(B_0)_P+\operatorname{trdeg}_{A/\mathfrak m}\kappa(P)=\dim A\). The Rees DVR has residue transcendence degree \(\dim(B_0)_P-1\) over \(\kappa(P)\). Transcendence degrees add to \(\dim A-1\). Its value group is discrete rank one and centre on A is the closed point. Thus it has precisely the divisorial definition in Rond--Spivakovsky. This establishes density on the actual all-support pointwise section without any singular SNC or characteristic-zero resolution argument.

Now assume A is normal. N2 gives, for each centred divisorial \(\nu\), a finite constant with \(\nu(f)\le C_\nu\operatorname{ord}_{\mathfrak m}(f)\) for nonzero nonunits; units are separately zero. Normalize \(\nu\) by \(\mathfrak a\). If \(\mathfrak m^r\subset\mathfrak a\), then every \(v_0\) in the section has \(v_0(\mathfrak m)\ge1/r\). Choosing a nonzero generator attaining a proper primary ideal's finite \(v_0\)-value yields \(\nu(I)\le C_\nu r v_0(I)\). For unit ideals it holds too. Therefore every finite real negative ideal combination, and its uniform limit \(g\in C_A\), satisfies
\[
g(\nu)\ge C_\nu r\,g(v_0).
\]
If \(g(v_0)=0\), nonpositivity forces g to vanish at all divisorial points; density and continuity force \(g=0\). If \(u=\inf_i g_i\in D_A\) has value zero anywhere, every nonpositive \(g_i\) is zero there and hence everywhere. Thus on a normal local branch every completed function is either identically zero or everywhere strictly negative, with bottom allowed. Compact upper semicontinuity bounds each negative branch above by a finite negative constant.

Let \(u\in D_S\) be normalization-fibre constant. Call its finitely many normal centres active if their functions are everywhere negative, including bottom centres. Active and inactive points cannot share a fibre. The tuple \(\sigma=-1\) on active centres, zero elsewhere, belongs to \(C_S\) by primary-gauge/unit gluing from n3 and is fibre constant. The independently proved bounded theorem n4 descends it WITH POSITIVITY to \(\sigma\in C_R\). Its active set \(T=\{\sigma=-1\}\) is clopen. The descent h of u equals zero off T and is at most \(-c\) on T for some finite \(c>0\). If T is empty the result is zero. If every active value is bottom, the decreasing sequence \(N\sigma\) already proves the result.

For the remaining cases, every \(u-\delta\), \(\delta>0\), is an everywhere negative fibre-constant \(D_S\)-function by n3's addition rule, so n5 gives \(h-\delta\in D_R\). We do NOT take their increasing limit. Instead take any finite continuous test \(F>h\) on compact T and a finite positive lower gap e. Choose \(0<\eta<1/2\) with \(\eta(\|F\|+e)<e/4\), and set \(\delta=\eta c/[2(1-\eta)]\). At finite h on T,
\[
(1-\eta)(h-\delta)=h+\eta(-h-c/2)>h.
\]
Since \(h\le F-e\), the same choice gives \((1-\eta)(h-\delta)<F-e/2\), also at bottom. A decreasing \(C_R\)-net \(g_i\) to \(h-\delta\) has one index satisfying \((1-\eta)g_i<F-e/4\) throughout T, by compact strict-test covers. Choose \(N>0\) with \(-N<\min_T F\). Then
\[
q=\max\{(1-\eta)g_i,N\sigma\}\in C_R
\]
is strictly between h and F on T and EXACTLY zero off T. Both terms are below F on T; the first dominates the strictly larger scaled target at finite h, and q is finite at bottom. Off T its second term is zero and its first nonpositive.

The set of these continuous majorants, exactly zero off T and strictly above h on T, contains zero and is downward directed by insertion with the continuous minimum of two tests. The Urysohn construction in n5, applied on T, makes its pointwise infimum exactly h. Reverse pointwise order therefore gives an actual decreasing \(C_R\)-net. This proves full normalization descent, including mixed finite/bottom and mixed-zero functions, without an increasing closure or unproved flattening.

For arbitrary \(R\hookrightarrow B\), pull back to \((B\otimes_R S)_{\rm red}\), injective over S by n4. N3's traces at each normal local centre descend the pulled-back fibre-constant function to \(D_S\); it is precisely h pulled back from R. Full normalization descent then gives \(h\in D_R\). Conversely ideal extension preserves uniform and decreasing limits. Every primary change follows n1. Empty zero-dimensional sections have just their unique empty function. Artifacts: n6.md and the paper. Rerun: python build\_delivery.py. These are mathematical draft proofs, not consequences of registration.

\section{n7. Discharge of the primary-source geometry and real perturbation}
\textit{Earlier local supports: n1.}

Use n1 for rings and the boundary. We prove the three geometric inputs, using the exact published primary units recorded below. No supplied continuation label is a premise.

Let Y be the actual blowup of a proper primary ideal I on a POSITIVE-DIMENSIONAL reduced Noetherian local R. Dimension zero is the empty-boundary case of n1, treated separately. No minimal component of this positive-dimensional local reduced ring is zero-dimensional: that would make its maximal ideal a minimal prime and force dimension zero. The Rees algebra is a subalgebra of R[t], so is reduced. Its component primes are obtained from the minimal primes of R: projecting the Rees algebra to each component domain gives prime kernels with intersection zero; component-separating elements of R show these kernels are the minimal primes. Since I has nonzero image on each positive-dimensional component, every component of its blowup dominates a component of R. The standard f-chart is the degree-zero subring R[I/f] of R\_f. Multiplication by f is injective there and I becomes (f). Consequently I O\_Y is invertible with a regular local generator, defining an ACTUAL effective Cartier divisor E; O\_Y(-E) is the relatively ample tautological bundle. The morphism is projective and is an isomorphism away from the closed point. Relative and absolute ampleness agree over this affine base, by complete Stacks 29.38.4--5 (01VG).

G1: T=H\textasciicircum{}0(Y,O\_Y) is a finite R-module by complete Stacks 30.16.1 and 30.16.3 (0B5S). Restriction to all generic components embeds T in the total quotient ring of R, since Y is reduced and those components are dense. The inclusion R into T is injective. Global sections localize on the base: choose a finite affine cover of the separated Y; its pairwise intersections are affine, and H\textasciicircum{}0 is the kernel of the difference map between finite products of the affine section modules. Localization is exact and commutes with those finite products. Since Y is the identity off the closed point, (T/R)\_p=0 for p different from m. The finite module T/R is therefore killed by m\textasciicircum{}h for some h>=0. Fix ONCE on this model
\[
\mathfrak c_0=\mathfrak m^{h+1},\qquad \mathfrak c_0 T\subset\mathfrak m\subset R.
\]
The containment is as elements of R: m\textasciicircum{}h T lies in R and then one more factor m puts it in m. This primary correction is not the full normalization conductor.

G2: Complete Stacks 28.27.13(7) (01PR) gives eventual generation of F tensor A\textasciicircum{}k for any coherent F and ample A on the Noetherian scheme Y. Thus for each invertible L, A\textasciicircum{}\{k-1\} tensor L is generated for large k. Complete 28.27.5 says ample tensor generated invertible is ample, so A\textasciicircum{}k tensor L is ample. Complete 28.27.2 gives the ample power-root assertion. These are applicable to Y without normality or smoothness. Their proofs were read, including the affine-section criterion and the high-power generation equivalences; deeper referenced scheme-theoretic results are published background, not claimed independently reconstructed.

G3 first for integral bundles: restrict to the full projective closed fibre X over the residue field, which may be nonreduced. Corrected Keeler, math/0108068v5 Corollary 5.19, gives generation of A\textasciicircum{}m tensor N uniformly over ALL integral curve-nef N for m sufficiently large (a fixed A has a very ample power). Apply it with N\textasciicircum{}k and m=k-1. Thus G=A\textasciicircum{}\{k-1\} tensor N\textasciicircum{}k is generated, and
\[
(A\otimes N)^k=A\otimes G
\]
is ample by 28.27.5. Taking a power root shows A tensor N ample on X. Complete Stacks 30.21.4 (0D2N) extends fibre ampleness to relative ampleness on a neighbourhood of the closed point. Every open neighbourhood of the closed point in Spec of a local ring is the whole base. The argument therefore proves integral curve-nef plus ample on Y. It applies directly to the nonreduced fibre; reduction is unnecessary. If used, complete 30.17.5 (01XO) separately supplies nilpotent invariance of ampleness.

The Keeler correction is essential. The complete immediate 1.5 proof, its 4.9--4.10 regularity/generation inputs, and the complete 5.19 derivation were read together with the ENTIRE appended 2018 erratum. False original 3.3 (geometric integrality without the needed change of ground field) is not used. E1.3 replaces it with a finite faithfully flat altered base after localization and smooth geometrically integral fibres over the new base. E2.4 transfers vanishing through that finite base; E2.5 repairs the induction proving 1.5. The erratum retains 1.5 and hence 5.19 on arbitrary projective schemes over arbitrary commutative Noetherian rings. No normal, reduced, characteristic, or smooth hypothesis is introduced. Its historical de Jong alteration, EGA, Illusie, Kleiman and regularity dependencies are disclosed published background. We do not claim every historical edition has been independently read.

Here is the real-Cartier extension, proved explicitly. Write additive bundle notation in \(\operatorname{Pic}(Y)\otimes_{\mathbb Z}\mathbb R\). Call a real class real-ample when it is a finite positive real sum of integral ample classes. Rational curve-nef F plus integral ample A is real-ample: for rational t>=0, clearing a common denominator makes A+tF genuinely rational ample by the integral theorem just proved. For arbitrary real t, bracket t by two nonnegative rational parameters and take their convex combination. Hence A+tF is real-ample. If S=sum s\_j A\_j is real-ample, with s\_j>0 and s=sum s\_j, then
\[
S+F=\sum_j s_j(A_j+F/s)
\]
is real-ample for rational nef F.

Fix an integral ample H and integral classes L\_i spanning a real curve-nef D. G2 gives positive integers N\_i such that N\_i H+L\_i and N\_i H-L\_i are ample. On every contracted integral curve, their positive degrees imply |deg L\_i|<=N\_i deg H. For rational delta>0 choose rational D\_r in this finite span with error coefficients e\_i in D-D\_r satisfying sum N\_i|e\_i|<delta/4. Then F=D\_r+delta H/4 is rational curve-nef. The residual is explicitly real-ample:
\[
D+\delta H-F=(3\delta/4-\sum_i N_i|e_i|)H+
\sum_i |e_i|(N_iH+\operatorname{sign}(e_i)L_i).
\]
The H coefficient is positive; omit zero error terms. Therefore D+delta H=F+S is real-ample. For arbitrary real delta>0 first use a rational delta\_0 between zero and delta and add (delta-delta\_0)H. This establishes curve-nef real plus any positive ample perturbation.

We also need finite-span openness, not a finite-rank assertion about the whole Picard group. For an integral ample A and finitely many L\_i choose N\_i A +/- L\_i ample. If sum N\_i|e\_i|<1, the identity
\[
A+\sum_i e_iL_i=(1-\sum_i N_i|e_i|)A+
\sum_i|e_i|(N_iA+\operatorname{sign}(e_i)L_i)
\]
is an ample-generator expression. For a general real-ample S reserve one positive summand sA, apply this expression to perturbations of that summand, and retain the other summands. This proves openness on every finite-dimensional coefficient span, even when relations exist among its classes.

Finally a rational Cartier class that is real-ample is genuinely Q-ample. Let G be the subgroup of Pic(Y) generated by its finitely many integral constituents and the integral ample classes in a positive real expression. G is a finitely generated abelian group, with finite torsion; \(G\otimes_{\mathbb Z}\mathbb R\) embeds in \(\operatorname{Pic}(Y)\otimes_{\mathbb Z}\mathbb R\) because the REAL NUMBERS are flat over Z. This tensor coefficient field is not the local base ring R. In G modulo torsion, equality with the rational class is a finite rational linear system for the positive expression coefficients. A positive real solution is approximated by a positive rational solution in that rational affine solution space (row reduction over Q proves rational points dense). Clear denominators and then kill the finite torsion. A positive integer multiple of the rational class is a tensor product of positive powers of integral ample bundles, hence ample. This proves genuine Q-ampleness without asserting finite rank of Pic(Y) or a Neron--Severi group.

Artifacts: n7.md and source-application-ledger.md. Rerun: python build\_delivery.py. These source applications discharge G1--G3 rather than assume unnamed positivity theorems. Their published recursive dependencies remain explicit; execution cannot certify them.

\section{n8. Independent geometric construction and unconditional comparison}
\textit{Earlier local supports: n1, n3, n4, n5, n6, n7.}

Use n1, n3 and n7 for the construction and comparison, and n4--n6 for its descent rules. Assume the local reduction has positive dimension; dimension zero has the unique empty-boundary cone directly. Here is the geometric definition BEFORE its comparison to ideals. Admissible models are actual blowups Y\_I of proper primary ideals of the reduced excellent local germ. On such a model consider an exceptional real Cartier divisor
\[
D=-\sum_{j=1}^r d_j A_j,\qquad d_j\ge0,
\]
where every A\_j is an ACTUAL effective integral Cartier divisor supported on the closed fibre. Real Cartier divisors here lie in the integral Cartier divisor group tensored over Z with the REAL NUMBERS; their bundle classes lie in the corresponding real Picard group. Require deg\_C D>=0 for EVERY proper integral curve contracted to the closed point. This is line-bundle degree positivity, and actual Cartier effectivity; it is not effectivity of a list of nonnormal Weil coefficients. Neither ideal generation nor membership in a completed function cone defines the positivity test.

The model system is directed. Y\_\{IJ\} maps to both Y\_I and Y\_J. To justify this over a reduced ring, on Y\_\{IJ\} the product of the pulled-back ideals is invertible, locally (f) with f regular. Set J'=f\textasciicircum{}\{-1\}J in the total quotient sheaf. I J'=O, which expresses 1 as a finite sum of products. In a local ring one product ij' is a unit. For any x in I, xj' belongs to O, and x=(xj')/(ij') times i. Thus I is principal there, with a regular generator because ij' is a unit in the total quotient ring. J is its invertible inverse times (f). The universal property of the blowup gives the two maps. Pullback preserves actual effectivity. For a contracted integral curve its image is a point or an integral curve; in the latter case degree of a pulled-back bundle is its image degree times a positive field degree, and in the former case it is zero. Thus curve-nefness is preserved. Identifying data after a common pullback is legitimate.

E\_I is a Cartier order unit on Y\_I, not merely a numerical order unit. An integral exceptional Cartier F is trivial as a divisor off the closed fibre. The coherent modules (O(F)+O)/O and (O(-F)+O)/O, viewed in the total quotient sheaf, are supported on that fibre. On a finite affine cover a common power of I O kills them. Since I O=O(-E), this gives inclusions O(F-NE) subset O and O(-F-NE) subset O. An invertible subsheaf of O with a regular local generator defines an actual effective Cartier divisor. Therefore
\[
-NE\le F\le NE
\]
in actual Cartier order, for some integer N. For finite real sums this yields real order bounds too.

Define evaluation on ALL boundary points as follows. For v with support p, first pass to the valued field Frac(R/p). A generator f of I of least finite value gives the centre on the f-chart, because every a/f has nonnegative value. The centre is independent of this choice by the blowup universal property, or by the overlap units when two generators have the same least value. Pull O(F) back to the spectrum of this field's valuation ring. The support's generic point lies outside the exceptional fibre, so O(F) there has its canonical trivialization as O. This identifies the pulled-back invertible module with a rank-one lattice inside the quotient fraction field. If g is a generator of that lattice, set f\_F(v)=-v(g). Its image is nonzero by the generic trivialization. This is a pullback of an invertible bundle with generic trivialization; it does NOT assume a homomorphism from the total quotient ring of R to Frac(R/p). Different generators differ by units of the valuation ring and so have value zero. Extend linearly to real divisors. Evaluation is homogeneous on rays, additive, invariant under common pullback, and f\_E=v(I). The order-unit bounds make it finite and bounded on a primary section. This construction uses quotient fraction fields; it never loses a nonminimal support by trying to use a fraction field of R itself.

We now prove continuity and comparison, rather than stipulate them. For a negative rational relatively ample exceptional L, take arbitrarily large sufficiently divisible k for which O(kL) is generated. Its sections F\_k=H\textasciicircum{}0(Y,O(kL)) lie in the finite algebra T of n7 because O(kL) subset O. They are finite R-modules by 0B5S. With the FIXED primary ideal c\_0 of n7 set J\_k=c\_0 F\_k. It is an ideal of R: multiplication by R preserves sections, and c\_0 T subset m AS ELEMENTS OF R. It is proper. Away from m it is the unit ideal, because c\_0 is unit there and O(kL)=O there; localization of H\textasciicircum{}0 and generation give this assertion. Therefore J\_k is proper primary. At every valued centre, global generation says the images of the sections generate the local invertible module, and at least one coefficient is a unit. Thus
\[
-v(J_k)/k=f_L(v)-v(\mathfrak c_0)/k
\]
on ALL semivaluations. All terms are finite. Since c\_0 is primary, its value is uniformly bounded. Hence f\_L is a uniform limit of negative rational primary profiles, with an explicit error bound ||v(c\_0)||/k.

For an arbitrary integral exceptional F, G2 and the order-unit assertion let us choose N so L=F-NE is both actual negative and ample. Apply the preceding section construction for any generated divisible k. Its equality gives
\[
f_F=-v(J_k)/k+v(\mathfrak c_0)/k+Nv(I).
\]
Each ideal in this equality is primary. This proves continuity of every exceptional Cartier evaluation, without assuming continuity of nonnormal model coefficients. Real evaluations are finite sums, hence continuous. A negative actual effective divisor has nonpositive evaluation.

Let M\_geom be all evaluations of the independent negative curve-nef divisors defined above, including zero; set C\_geom to their uniform closure, and D\_geom to actual infima of nonempty decreasing C\_geom-nets, including bottom. We prove M\_geom subset C\_R. Write D=-sum d\_j A\_j. By n7, D-delta E is real-ample for every delta>0, since -E is ample. In the finite span of the A\_j and E, round ALL nonnegative coefficients UP to nearby rational numbers. Finite-span openness retains ampleness; the rounded divisor D' is actual negative and rational. It is genuinely Q-ample by the finite-subgroup argument in n7. Continuity and boundedness of the finitely many evaluations give ||f\_D'-f\_D|| arbitrarily small as delta and rounding errors tend to zero. Apply the fixed-c\_0 section construction to D' and send divisible k to infinity. This proves f\_D in C\_R. The generated section profiles can be chosen below f\_D, but only uniform approximation is needed. No full-conductor value is divided by k.

Conversely every raw finite REAL negative ideal combination is geometric: pull the finitely many primary-ideal exceptional divisors to their common product blowup. Each -E\_i is a pullback of an ample bundle and is curve-nef; each E\_i is actually effective. The combination meets the definition and evaluates to the given profile. Unit terms evaluate to zero. Thus
\[
Q_R\subset M_{\rm geom}\subset C_R,\qquad
C_{\rm geom}=C_R,\qquad D_{\rm geom}=D_R.
\]
This is the additional geometric comparison theorem required by ORIGINAL P12.2. It is not a renaming of ideal positivity. Raw equality with Q\_R is not claimed.

The geometric meaning is relative Cartier degree positivity on the explicitly stated primary-model system, followed by uniform and decreasing completion. We make no cofinality assertion for arbitrary proper birational models, no singular SNC or positive-characteristic resolution assertion, and no replacement of actual effectivity by nonnormal Weil coefficients. Finite pullback can also be checked on cores: a base primary ideal extends to a boundary ideal on the finite semilocal source; its blowup maps to the base blowup, actual effectivity pulls back, and degrees have the point/positive-degree rule. The comparison and n1, n4--n6 prove the full gauge and conductor rules without needing a separate raw descent claim.

Artifacts: n8.md and P12.2-paper.tex/PDF. Rerun: python build\_delivery.py. The independent curve test is the new geometric content; the uniform comparison is a proof. No envelope, product, energy, CLD, or classical radial-extension identification is asserted.

\section{n9. Full smooth BFJ specialization on all semivaluations}
\textit{Earlier local supports: n1, n3, n8.}

Use n1, n3 and n8. The published BFJ application in this note is ONLY at the closed point of a smooth affine variety over an algebraically closed field of characteristic zero. Dimension at least two is treated by BFJ; dimensions one and zero are treated directly below. No characteristic-free singular resolution is inferred.

The complete used units of Boucksom--Favre--Jonsson, Valuations and plurisubharmonic singularities, were checked from the bound publication: the setting; semivaluations and primary profiles; simple modifications, simplices, and retractions in Sections 1.2--1.7; Definitions 2.1 and 2.5, Proposition 2.2, and 2.7--2.8; 3.6--3.14; and the complete proofs of 3.9--3.10 through 3.18 and 3.22--3.25. The historical smooth resolution, usual multiplier ideal subadditivity, and smooth Izumi inputs remain disclosed published background. BFJ's nef Weil cone is its weak closed cone generated by formal-ideal divisors; its formal psh functions are finite and strongly continuous on quasi-monomials and have that nef coefficient divisor. They extend to ALL semivaluations by
\[
H_W(v)=\inf_\pi H_W(r_\pi v).
\]
Proposition 2.7 says H\_W<=H\_W composed with r\_pi; the composed functions decrease under refinement. The projective-limit statement 1.13 applies to the all-support compact semivaluation space, not just fraction-field valuations.

Proposition 3.11 proves substantive continuity on each CLOSED simplex. Its proof uses multiplier profiles h\_k=(1/k)log|L\textasciicircum{}2(kW)| and the estimate |h\_k-H\_W|<=A/k on divisorial points. Thinness A is bounded on each closed simplex. The continuous multiplier profiles are consequently uniformly Cauchy there, and their limit supplies the continuous extension, including its faces. One must not substitute the false assertion that a finite convex function is automatically continuous at the boundary of a closed simplex.

Correct 3.13 as follows. For a nonzero formal psh g, define, for integers j>=1,
\[
B_j=L^2(2^j g)+\widehat{\mathfrak m}^{\,4^j},\qquad
g_j=2^{-j}\log|B_j|=\max\{2^{-j}\log|L^2(2^jg)|,-2^j\}.
\]
The exponents were checked visually in the actual PDF on printed page 475. Subadditivity 3.10 gives the decreasing first term, and 3.9 gives convergence to g at every quasi-monomial. The truncation also decreases. The ideals contain a boundary power but early terms can be the unit ideal. Nonzero g has g(ord\_m)<0 by 3.14, and 3.9 at ord\_m makes the multiplier profile strictly negative there for all large j. Thus B\_j is EVENTUALLY proper primary; discard a finite prefix. For g=0 use the UNIT ideal, not a sequence of proper primary ideals decreasing to zero. The identically bottom function is not a finite formal psh function on quasi-monomials; add it separately by -N=-v(m\textasciicircum{}N).

The approximants decrease to g on ALL semivaluations. Given finite t>g(v), choose a retraction with g(r\_pi v)<t by the definition of the all-retraction infimum, even when g(v) is bottom. Then choose j with g\_j(r\_pi v)<t. The formal-psh inequalities give
\[
g(v)\le g_j(v)\le g_j(r_\pi v)<t.
\]
This proves full convergence; it is not merely divisorial or quasi-monomial convergence. N1 contracts each completed primary B\_j exactly to R, preserving the profile. The completed unit also contracts to the unit. Thus every BFJ function, with bottom adjoined, belongs to D\_R, and every finite continuous BFJ function belongs to C\_R by the compact decreasing-continuous uniform-convergence argument (Dini's open-cover proof).

Conversely any finite real primary profile is BFJ formal psh, since its coefficient divisor is nef and its evaluation is continuous on simplices; principalization and the published retractions give its canonical all-semivaluation extension. Let g be a uniform limit of such profiles. Its divisorial coefficients give a finite weak limit W in BFJ's closed nef cone. Proposition 3.11 gives H\_W. Uniform convergence on quasi-monomials shows g=H\_W there: they agree on dense divisorial points of each simplex and both are continuous there. Moreover the same uniform error estimate survives taking the infimum over ALL retractions. Thus g(v)=inf\_pi H\_W(r\_pi v) at every semivaluation. This proves C\_R is EXACTLY the finite continuous BFJ subclass.

Now let g=inf\_i g\_i be a decreasing C\_R-net. BFJ 3.14 gives g\_i(v)<=g\_i(ord\_m) at every v (first on quasi-monomials, then by retractions). If g(ord\_m) is bottom, g is bottom everywhere. Otherwise 3.14 also gives C(v)g\_i(ord\_m)<=g\_i(v) at every quasi-monomial v, so g is finite there. Its divisorial coefficients give a finite weak-nef W. On a simplex the decreasing limit g is finite convex. For any point x take its MINIMAL face; x lies in that face's relative interior, where finite convex functions are continuous. H\_W is continuous on the whole face by the actual 3.11 theorem. Their equality on the dense rational divisorial points of that relative interior proves equality at x. Vertices are already divisorial. This handles every face without asserting closed-simplex continuity from convexity alone.

Finally, extended infima commute over the Cartesian product of the net and ALL models:
\[
g(v)=\inf_i\inf_\pi g_i(r_\pi v)
=\inf_\pi\inf_i g_i(r_\pi v)
=\inf_\pi H_W(r_\pi v).
\]
No arbitrary decreasing-net flattening principle is used. Hence D\_R minus the identically bottom function is exactly the full BFJ formal psh class with its canonical extension to all semivaluations. D\_R is that class with bottom adjoined, and n8 gives the same statements for D\_geom and C\_geom.

For a smooth DVR every nonempty normalized boundary is a singleton. Primary ideals are powers of the parameter, so C is all finite nonpositive constants and D all extended nonpositive constants. The closed point itself is an effective Cartier divisor, and degree conditions on contracted curves are empty, giving the same independent geometric core. In dimension zero the centred boundary is empty and all statements mean its unique empty function. General primary gauges in the BFJ range follow n1's exact homogeneous cocycle, with bounded positive denominators preserving uniform and decreasing convergence. Completion homeomorphism and primary contraction apply without claiming equivalent model categories.

Artifacts: n9.md, reports/source-images/bfj-475.png, source-application-ledger.md. Rerun: python inspect\_primary.py; python build\_delivery.py. This specialization has precisely the published smooth affine algebraically closed characteristic-zero domain; it does not transfer singular SNC theory to arbitrary excellent germs.

\section{n10. Nonisolated cusp: separation from the entire finite-real raw cone}
\textit{Earlier local supports: n1, n3, n4, n8.}

Use n1, n3, n4 and n8. Work in characteristic zero, with
\[
R=k[t^2,t^3,z]_{(t^2,t^3,z)},\quad S=k[t,z]_{(t,z)},\quad
\mathfrak a=(t^2,t^3,z),\quad \mathfrak c=t^2S.
\]
Indeed S is finite via t, is normal, and has the same fraction field. R consists of polynomials with no term of t-degree one, localized at the stated maximal ideal; multiplying by t shows its conductor is exactly t\textasciicircum{}2S. The quotient R/c is k[z] localized, and S/c is its square-zero t-thickening. Their all-semivaluation boundaries are the same singleton when nonempty. N1 therefore makes restriction V\_\{aS\}(S) to V\_a(R) a continuous bijection and hence a homeomorphism, including the support (t) upstairs. This is a two-dimensional nonisolated cusp, not just a finite curve link.

Write s=w(t), u=w(z), b=z+t and B=w(b); min(2s,u)=1. Allow s,u,B to be infinite whenever the normalization permits. The bounded continuous positive profile
\[
A(w)=w((b,t^2))=\min\{B,2s\}
\]
descends topologically. For k>=1 define actual downstairs ideals
\[
K_k=(b^k(z-kt),\ t^2b^k,\ t^{2k+2})\subset R.
\]
The first generator has zero t-linear coefficient: the two contributions k z\textasciicircum{}k and -k z\textasciicircum{}k cancel. The other two have t-degree at least two. Localizations do not change membership. The ideals are proper primary: the last generator puts t\textasciicircum{}2 and t\textasciicircum{}3 in the radical, and the first then puts z in it.

For EVERY normalized semivaluation, every generator has value at least kA. For the first this is kB+w(z-kt)>=kA, and the other two are immediate. The upper bound is a complete case analysis. If A=2s<=B, cancellation forces s=u=1 (different values would give B=min(s,u)<2s), and the last generator has value kA+2. If A=B<2s and s<=u, normalization gives s<=1, and the second generator has value kA+2s<=kA+2. If s>u, including s infinite, then u=1 and B=w(z-kt)=1; the first generator has value kA+1. The case u infinite belongs to the second case. Thus on ALL supports, including cancellation and infinite coordinates,
\[
0\le v(K_k)/k-A\le2/k.
\]
Hence -A lies in C\_R, and is a normalization-raw profile in N\_R. This analytic bound is a proof; a finite sample check is only a diagnostic.

We separate A from the ENTIRE finite REAL positive raw cone P\_R, not merely rational single-ideal profiles. Let w\_infty kill t and give z value one. Let w\_star be the monomial valuation in coordinates (t,b) with weights (1,2); then z=b-t has value one, so both valuations are in the same primary section. They give A(w\_infty)=1 and A(w\_star)=2.

For ANY proper primary downstairs ideal I, put d=w\_infty(I), a positive integer. Choose f in I attaining d; its local denominator is a unit at both valuations. Expand in k[[t,z]]: no coefficient of t-degree one occurs. Its degree-d homogeneous part has nonzero z\textasciicircum{}d coefficient and zero t z\textasciicircum{}\{d-1\} coefficient. In characteristic zero it cannot be a scalar multiple of b\textasciicircum{}d=(z+t)\textasciicircum{}d. After substituting z=b-t, it therefore contains a nonzero term t\textasciicircum{}\{d-j\}b\textasciicircum{}j with j<d. That monomial has w\_star-value d+j<=2d-1. Terms of other total degrees cannot cancel it, since the linear substitution preserves total degree. If other terms have still smaller value they only strengthen the bound. Consequently
\[
w_\star(I)<2w_\infty(I).
\]
Every nonzero finite real positive combination of proper primary gauges retains this strict inequality; unit terms are zero. But A has equality. Thus A is outside P\_R, and -A is outside Q\_R. We have proved N\_R is not Q\_R and C\_R is not Q\_R in this example. Their uniform closures still agree by n4. This does not assert raw geometric and ideal equality.

Finally, v(c)<infinity is an OPEN complement of the CLOSED infinity stratum, not generally closed. Monomial valuations w\_N(t)=N, w\_N(z)=1 converge pointwise to w\_infty. Their conductor values are 2N, while w\_infty(c)=infinity. The quantities v(c)/N are therefore not uniform errors on the whole section. Neither descent nor comparison uses such an error globally; n4 uses quotient estimates at infinity, and n8 uses a fixed PRIMARY correction.

Artifacts: n10.md, exact\_checks.py, reports/exact-checks.json. Rerun: python exact\_checks.py; python build\_delivery.py. The script reports its exact finite input domain and does not replace the universal proof or the raw-cone separation argument.

\section{n11. Rejected shortcuts and the radial scope constraint}
\textit{Earlier local supports: n1, n2, n3, n4, n8, n10.}

Use n1--n4, n8 and n10. These are diagnostics, not extra hypotheses.

An incompatible-contraction shortcut is false even over a smooth normal surface. Inject A=k[x,y]\_(x,y) diagonally into A times A and take J=((x,y\textasciicircum{}2),(x\textasciicircum{}2,y)). Then I\_N=J\textasciicircum{}N intersect A=(x,y\textasciicircum{}2)\textasciicircum{}N intersect (x\textasciicircum{}2,y)\textasciicircum{}N. The m-order of both fibre ideals is one. A monomial x\textasciicircum{}a y\textasciicircum{}b belongs to I\_N exactly when 2a+b>=2N and a+2b>=2N. Summing gives a+b>=ceil(4N/3); balanced nonnegative exponents with that sum achieve both inequalities. Therefore
\[
\operatorname{ord}_{\mathfrak m}(I_N)=\lceil4N/3\rceil,
\]
whose normalized limit is 4/3, not one. The two ideal profiles were not fibre constant. No general contraction-as-fibre-minimum formula is used in n4--n6.

The literal positive-gauge inequality in BdFFU Proposition 2.4(b) has the wrong sign. At a smooth surface take parameters x and b=x+y, and the monomial valuation v(x)=1, v(b)=2. On the first-blowup dual complex its retraction is ord\_m. For I=(x+y,x\textasciicircum{}3), v(I)=2 whereas r(v)(I)=1. Hence the printed inequality r(v)(I)>=v(I) fails. The negative function -v(I) has the expected opposite inequality. Complete Definition 2.1--2.2 and Proposition 2.4 with proof were read from the bound BdFFU paper, solely for this diagnostic. Its own normal complex-variety hypotheses do not repair the sign. It supplies no premise of the construction.

The geometric boundary cone is not asserted to be the cone of literal CLD or classical logarithmic radial extensions. On a nodal curve R=k[x,y]\_(x,y)/(xy), its finite normalization has two DVR centres. The tuple (-1,0) is in C\_S and is fibre compatible: the conductor has finite positive value on both boundary points, so its restriction has no identified pair. N4 therefore puts (-1,0) in C\_R, and n8 puts it in C\_geom. For a classical analytic node with two normalized discs, a literal zero logarithmic slope on one disc forces the nodal value to be at least zero by upper semicontinuity. The other literal negative slope tends to minus infinity at its origin. If its nodal value is at least zero, the subharmonic maximum principle on that disc (the outer-circle value is zero and interior values are negative) forbids the resulting interior maximum. Thus this mixed-zero profile cannot be a literal classical psh radial extension. The original dossier's corresponding CLD closed-disc obstruction is kept as a scope constraint, not silently overridden or used as a theorem premise. We make no CLD comparison theorem for other singular germs. A boundary trace of a different kind is a separate question.

Artifacts: n11.md, exact\_checks.py, reports/exact-checks.json. Rerun: python exact\_checks.py; python build\_delivery.py. The monomial computation is proved for all N; finite enumeration only checks a specified range. The sign example is exact. This note rejects invalid approaches while retaining the independent geometric construction and its precise scope.

\section{n12. Complete ORIGINAL P12.2 candidate and deliverable evidence}
\textit{Earlier local supports: n1, n2, n3, n4, n5, n6, n7, n8, n9, n10, n11.}

Use every earlier local note n1--n11. We claim a complete solution of the exact ORIGINAL P12.2 target, under the assignment's permitted published-background rule, as an unverified mathematical draft requiring subsequent independent whole-original review.

For every excellent singular or nonnormal local germ, reduce nilpotents and take its intrinsic centred ray boundary. Actual primary blowups carry an independent geometric core: negative ACTUAL effective finite real Cartier presentations having nonnegative degree on every contracted proper integral curve. Evaluation uses the support quotient fraction field and its valued centre. The admissible model system, common pullbacks, order units and nonminimal supports are explicit. N7 proves the exact primary-source G1--G3 applications, including the corrected Keeler v5 route on arbitrary nonreduced projective fibres and every real perturbation/rationalization transition. N8 proves C\_geom=C\_ideal and D\_geom=D\_ideal; degree positivity was defined before that comparison, so this is not merely an ideal completion renamed.

N1 proves the intrinsic primary-gauge cocycle, completion and exact normalization conductor fibres. N2--n6 prove full continuous-uniform and actual decreasing-net descent for EVERY reduced excellent finite injection, including unbounded, allzero, allbottom, partialbottom, mixedzero, conductor-infinite and nonminimal supports. Norm traces have generic inseparable lengths, quotient-support Gauss values and proved full fibre support. Conductor induction covers all minimal/zero-dimensional components and uses a positive boundary factor, a common graded exponent, the exact residue fibre product and compact open covers. Internal increasing profiles never serve as an assumed completion operation; actual completed descent is obtained by insertion and downward-directed continuous majorants.

N9 proves the FULL smooth BFJ specialization on ALL semivaluations in its actual smooth affine algebraically closed characteristic-zero domain. It uses substantive closed-simplex continuity, minimal-face interiors and commuting all-retraction infima, corrects zero/unit and eventual proper-primary approximants, and treats bottom, DVR and empty dimension zero separately. N10 separates the entire finite REAL raw cone from its normalization-raw and continuous completions by a nonisolated cusp, with the universal error <=2/k. N11 rejects the incompatible-contraction shortcut and wrong BdFFU sign and preserves the mixedzero radial obstruction. P12.3--P12.6 envelopes/products/energy/operator theory remain separate targets, not silently supplied.

No exact P12.2 mathematical gap is retained. The limits are explicit: only the stated primary-model system, no arbitrary-model cofinality, no nonnormal Weil substitute, no singular positive-characteristic SNC transfer, no literal CLD/classical radial identity, and no intersection-multiplicity theory for nilpotents. Published recursive historical dependencies are disclosed rather than claimed fully independently reread. There are zero accepted incoming notes and no provenance labels in support lists. The complete original-task.json, entire family dossier, all thirteen whole mathematical DATA files, and all complete immediate used primary units were read, with proof-body display truncations repaired. Hashes and old flags grant no mathematical status.

Fresh deliverables are P12.2-paper.tex and P12.2-paper.pdf, argument.md, source-application-ledger.md, exact-scope.md, whole-proof-coverage.md, original-scope-decision.md, native-notes.json, n1.md through n12.md, and the fresh scripts/inputs/reports. Native-notes contains the complete proof-bearing local chain, not just an index. Reports retain input and own SHA-256 hashes, environment details, local source-image inspection, exact finite checks, every PDF build attempt/warning, execution history, and final integrity/support audit. No immutable input was edited and no supplied builder/checker/binder was executed or imported.

The observed finite diagnostic result is PASS on 1600 cusp value cases (k=1..32, sixteen pairs of rational weights, coherent equal-weight cancellation cases and two infinite-coordinate endpoints), 32 exact t-linear cancellations, and all contraction minima for N=1..36 in the box 0<=a,b<=2N. The largest sampled value J-kA is exactly 2. These observations do not establish the universal mathematical claims; their proofs are n10--n11. PDF compilation and support/hash auditing likewise establish artifact consistency, not mathematical correctness.

Rerun from this directory: python audit\_inputs.py; python inspect\_primary.py; python exact\_checks.py; python build\_delivery.py; python final\_audit.py. The build uses pdflatex with -no-shell-escape. The source inspection only rasterizes three explicitly bound local pages. No network, provider, subagent or heavy compute was used; no supervised heavy-compute facility was provided or needed. Account/auth/runtime/HOME/CODEX\_HOME/proxy/CA state was not mutated. All fresh writes are in this workspace. The requested model designation is a caller execution request, not an independently measured theorem or provider attestation.

Candidate=true here means that the last note claims the entire ORIGINAL P12.2 construction, comparison, smooth specialization and gauge/conductor rules. It does not mean paper acceptance, catalogue closure or independent verification. Subsequent independent whole-original review is required.

\section*{Reproducibility and review status}
The accompanying source-application-ledger, exact-scope, whole-proof-coverage,
and original-scope-decision files state the exact domains, complete immediate
source units checked, recursive dependencies, limits and proof coverage.
The full native-notes JSON repeats every proof body and its earlier local supports.
Input and own hashes, finite diagnostic inputs/output, environment and every
build attempt are retained in reports. No supplied executable or network was used.
Execution establishes reproducibility only. The candidate flag is a draft
solution claim and does not accept the paper or close a catalogue.
\end{document}
\end{Verbatim}
\section{Mandatory source and scope annotation, unchanged}
\begin{Verbatim}[breaklines,fontsize=\tiny]
{
  "recorded_utc": "2026-10-08T12:26:53.361733+00:00",
  "atom": "P12.2",
  "continuation": 331,
  "source_unchanged": true,
  "original_statement_unchanged": true,
  "independent_geometry": "Actual primary blowups; actual negative effective finite REAL Cartier presentations; nonnegative degree on EVERY contracted proper integral curve, defined before ideal completion. Uniform comparison and actual decreasing-net comparison then proved.",
  "exact_general_scope": "Every excellent reduced singular/nonnormal local ring in positive dimension, arbitrary characteristic and residue field; nilpotents invisible, zero-dimensional local reduction boundary empty; ALL semivaluation supports retained.",
  "model_system_scope": "Directed actual primary blowups with pullback and arbitrary-support evaluations. Cofinality among ALL proper birational models and a full arbitrary-model b-divisor classification are unproved and unclaimed.",
  "smooth_specialization_scope": "Published BFJ smooth affine CLOSED germs over an algebraically closed characteristic-zero field, dimension>=2, plus direct DVR/empty cases. ALL-retraction extension, closed-face continuity, eventual proper-primary approximation, zero/unit and bottom exceptions retained. No singular/positive-characteristic SNC transfer.",
  "normalization_and_finite_scope": "Full conductor quotient; finite OPEN singleton stratum and CLOSED infinite-value quotient stratum; positive C and actual decreasing-net D descent, including mixed zero/partial bottom. Full finite-source fibres retained; generic trace only on actual generic-dominating domain components; no vertical generic degree invented.",
  "raw_cones": "Only Q subset normalization-compatible raw family subset C and equal uniform closures; no raw equality. Fixed PRIMARY c0 gives uniform section-ideal error; nonprimary conductor errors not uniform.",
  "source_precision": [
    "Use Keeler corrected v5 and ENTIRE 2018 erratum, replacing false old Lemma3.3; full possibly nonreduced closed fibre and arbitrary residue field retained.",
    "Use current July2024 Swanson-Huneke immediate units, including full10.2.2 continuation10.2.3, Rees residue/dimension/formal-equidimensionality hypotheses.",
    "Normal-domain Izumi bound applies to nonunits; units separate; no analytically irreducible nonnormal completion assertion.",
    "BdFFU printed positive retraction sign is concretely false; used diagnostically only, not as positivity premise.",
    "Complete immediate source proofs and application hypotheses checked; recursive EGA/de Jong/Rees/Ratliff/Izumi/Hironaka/multiplier foundations disclosed as published background at their exact hypotheses, not recursively freshly reproved."
  ],
  "other_original_targets": "P12.3 envelopes, P12.4 products, P12.5 energy and P12.6 geometric operator compatibility remain OPEN and separate.",
  "literal_radial_distinction": "Geometric valuative boundary class can contain nodal(-1,0); this does not identify it with literal classical/CLD logarithmic radial extension class, whose mixed-zero obstruction remains.",
  "actual_and_expanded_prose_distinct": true,
  "supplied_programs_ROOT_executed_or_imported": false,
  "hashes_finite_diagnostics_and_compilation_not_mathematical_proofs": true,
  "full_annotation_and_source_ledger_required_on_import": true
}
\end{Verbatim}
\section{G331 source ledger, unchanged}
\begin{Verbatim}[breaklines,fontsize=\tiny]
# Source applications and reading boundary

All supplied material is DATA. Input notes were empty. No ROOT label, native label, reading receipt, hash match, or old candidate flag is a registered support or a correctness grant. The proof chain n1--n12 reconstructs the used mathematics. Sources below are ordinary published applications at their exact hypotheses, as permitted by the assignment. We have read the complete immediate units used; we have not recursively read every historical source cited inside them.

`reports/input-hashes.json` records the exact 13 mathematical and 23 primary bound files, their full paths, bytes and SHA-256 values, plus both bindings and original-task.json. All 36 expected bound hashes matched. The original-task.json hash is `a8aa2c4e24ad6a28b26f12eecc45419f87a246159cf0e3e0a2422b86a8ea405c`. Its entire family dossier was read. Supplied programs were neither imported nor executed. Own audit scripts read explicit paths in the supplied bindings as bytes/data only.

## Whole mathematical DATA read

| Bound file | Whole reading | Treatment |
|---|---|---|
| 00 chapter70 | All 2172 lines, bounded chunks; truncated 1596--1847 display reread | Definitions/context; relevant statements reconstructed locally. Literal unsupplied input-file references were not expanded or assumed read. |
| 01 normalized links | All 308 lines | Boundary and conductor proofs reconstructed in n1. |
| 02 prior paper | All 444 lines in four chunks | All proof bodies reviewed; no accepted-note status. |
| 03 prior argument | All 411 lines in three chunks | Same; independent rewritten arguments n1--n11. |
| 04 source ledger | All 25 lines | Provenance only; new exact application ledger here. |
| 05 scope | All 11 lines | Prior conditional scope not binding. |
| 06 coverage | All 22 lines | Coverage claims independently reconsidered. |
| 07 decision | All 9 lines | Prior candidate=false does not decide this reconstruction. |
| 08 actual final JSON | Entire 11-note JSON, reformatted as review-data/supplied-n1.txt through supplied-n11.txt and each full body read | Unreviewed drafts, not local registered supports. |
| 09 reading binding | All 33 lines | Metadata only. |
| 10 acquisition metadata | All 152 lines | Metadata only. |
| 11 normal trace addendum | All 65 lines | Coefficient trace/full-support/finite-gap reasoning reconstructed in n3 and n5--n6. |
| 12 geometry working proof | All 80 lines | Every real perturbation transition checked and reconstructed in n7--n8; not accepted wholesale. |

## Immediate published primary applications

Line numbers refer to the exact bound text extraction, except the single HTML-only lemma. PDFs and their text counterparts are both hashed; this is not a claim to have visually read every PDF page.

| Source and complete units read | Exact application | Restrictions and recursive dependencies |
|---|---|---|
| Stacks 0B5S, 30.16.1 all five assertions/proof and 30.16.3, whole text 1--67 | Finite sections on proper schemes with an ample bundle over Noetherian base; finite T and F_k on primary blowups in n7--n8 | Usual Noetherian projective coherent cohomology is published background. H0 localization is separately proved by the finite affine-cover kernel. |
| Stacks 01PR, 28.27.2 and 28.27.5 (1--47), 28.27.13 all nine equivalences/proof (70--116) | Ample power roots, ample times generated bundle, and high-power generation; n7 G2, continuity shift, genuine Q-ample comparison | Affine-section ampleness criterion and preceding referenced scheme results disclosed. Not a singular resolution theorem. |
| Stacks 01VG, whole text 1--104, especially 29.38.4--5 | Relative ampleness equals absolute ampleness over affine local base; n7 | Relative ampleness definition and base locality used at exact projective scheme. |
| Stacks 0D2N, whole text 1--42, full 30.21.4 proof | Ample restriction to the possibly nonreduced closed fibre implies relative ample near the closed point, hence on entire local base | Immediate proof uses lifting from 30.21.3/formal functions, finiteness from 30.21.2 and finite pullback ampleness 29.38.7. Those recursive source bodies are not all separately read. |
| Stacks 01XO, full 30.17.5 statement/proof, 55--72 | Optional nilpotent invariance of fibre ampleness | Main Keeler route uses full nonreduced fibre directly, so no reduction is needed. |
| Keeler math/0108068v5: setting and 1.5 at 50--135; definitions 2.8--2.9, 2.13--2.14 at 210--286; complete 4.9--4.10 at 824--888; complete original 1.5 proof and 5.19 derivation at 1220--1288; ENTIRE appended 2018 erratum, 1780--2041 | 5.19: fixed ample powers uniformly generate twists by ALL curve-nef integral bundles on any projective scheme over any commutative Noetherian ring; yields integral ample plus nef on full fibre, n7 | Original 3.3 is false and expressly NOT applied. E1.3 changes to finite faithfully flat base with smooth geometrically integral fibres after localization. E2.4 and E2.5 repair 1.5; 5.19 remains valid. Historical de Jong, EGA, Illusie, Kleiman and regularity results are disclosed, not independently recursively reread. No char/smooth/normal/reduced hypothesis is added. |
| Swanson--Huneke CURRENT July 2024 author PDF: full 1.1.5 (580--624); 1.2.3--5 (650--722); 6.4.1--2 (5920--5979); 6.8.1--4 (6540--6630) | Support-quotient lifting agrees with n1's Zorn proof; all-minimal-component integral dependence criterion in n2, used for Q=K a^L in n4--n5; common exponent also proved by finite graded module argument | Zero-dimensional nonunit components must have Q-image zero; positive boundary factor L>0 ensures this. Book text-extraction bars repaired by visual page inspection. No assertion of byte identity with 2006 print. |
| Same book: complete 10.2.2 proof ending in 10.2.3 (9230--9390); 10.4.3 proof (9698--9756); B.3.2 and immediate context (19258--19310); B.5.1/proof (19435--19468) | Rees DVRs of positive-dimensional excellent local domains, residue trdeg under formal equidimensionality, exact birational dimension formula, universal catenarity to formal equidimensionality; n2 and density in n6 | B.3.2 describes a modification of its earlier dimension argument. Mori--Nagata, Ratliff, Rees historical dependencies disclosed. Zero-dimensional components handled separately. |
| Stacks 0C23 normal completion, full immediate HTML lemma 15.53.6 and proof | Normal excellent local domain has normal/domain completion and hence is analytically irreducible/unramified; n2 | Uses normal formal fibres and Algebra 10.163.8; latter not independently read here. No assertion for arbitrary nonnormal domain. |
| Rond--Spivakovsky 1404.4885v1 pages 1--2, text 1--115: Definitions 1.1--1.2 and Theorems 1.3, 1.4, 1.7 | Izumi--Rees comparability on an analytically irreducible local DOMAIN; asymptotic m-order and order error yield nonunit divisorial bound, n2 and zero dichotomy n6 | Units handled separately; literal strict Remark 1.6 not used. Original Rees 1989 and historical Izumi not independently read. |
| Swanson Rees-valuations survey, ENTIRE Section 5, text 833--921 | Check 5.2 really has complete-domain hypotheses | General analytic-domain application uses exact Rond--Spivakovsky 1.3 instead. |
| BFJ published Valuations and psh Singularities: complete setting and relevant Sections 1.2--1.7, 2.1--2.2, 3.1--3.6, text 1--1398 in bounded chunks; proofs through 3.25 | Full smooth BFJ recovery, n9, including substantive closed-simplex 3.11, corrected proper-primary 3.13 and all-retraction infima | ONLY smooth affine closed germ over algebraically closed characteristic zero, dimension>=2; DVR and empty dimension zero direct. Historical Hironaka, multiplier subadditivity, usual smooth Izumi disclosed. No arbitrary-characteristic singular SNC transfer. PDF page 475 inspected to settle 2^j and 4^j exponents. |
| BdFFU published singular multiplier ideals: complete Definitions 2.1--2.2 and Proposition 2.4 proof, text 175--303 | DIAGNOSTIC ONLY: positive-gauge 2.4(b) sign fails in explicit smooth example, n11 | No premise imported. Normal complex-variety hypotheses and literal sign issue preserved. |

The five bound geometry HTML counterparts are hashed data counterparts, not separately claimed line-by-line HTML proofs. The geometry packet has five Stacks HTML/text pairs plus Keeler PDF/text; the algebra packet has three PDF/text pairs plus one HTML lemma; the smooth packet has two PDF/text pairs, totalling 23 primary files. Immediate proofs were read from their text form, with three selected PDF pages inspected for extraction ambiguities.

## New proof boundary

Norm-coefficient fibre support, quotient Gauss identities, real uniform lattice closure, common graded exponent, bounded conductor approximation, insertion and mixed-zero descent, actual model degree positivity, the fixed-primary section correction, real perturbation and finite-Picard-subgroup rationalization are proved in the local note chain. They are not accepted on the strength of familiar names or supplied reading metadata. Small exact checks test only stated finite domains. Build success, support validation and matching hashes are reproducibility evidence, not mathematical acceptance.

Rerun locally: `python audit_inputs.py`; `python inspect_primary.py`; `python exact_checks.py`; `python build_delivery.py`; `python final_audit.py`. No network or supplied executable is needed.
\end{Verbatim}
\section{G331 exact scope, unchanged}
\begin{Verbatim}[breaklines,fontsize=\tiny]
# Exact scope

The construction is for every excellent local germ, after reduction, with a nonempty intrinsic centred ray boundary when the local reduction has positive dimension. Sections use arbitrary proper primary ideals. Nonminimal prime supports and infinite function values are retained. Zero-dimensional local reduction gives the empty boundary. For nonreduced germs this is a semivaluation theory of the reduction; nilpotent intersection multiplicities are not supplied.

The independent geometric core uses actual primary blowups of the reduced germ, actual negative effective finite real Cartier presentations, and degree >=0 on every contracted proper integral curve. Common product-ideal blowups identify presentations. There is no assertion that this system is cofinal in all proper birational models, no nonnormal Weil-coefficient substitute for Cartier effectivity, and no arbitrary-characteristic resolution assertion. The uniform and actual decreasing completions are justified by the proved comparison C_geom=C_ideal and D_geom=D_ideal; core positivity is geometric before that comparison.

P is the finite REAL positive ideal cone; Q=-P; C is its finite continuous uniform closure; D consists of actual decreasing C-nets with bottom allowed. N is the normalization-RAW compatible cone. Q subset N subset C, and both uniform closures equal C. The nonisolated cusp separates N and C from Q. No raw equality is asserted.

Primary-gauge changes are the exact homogeneous cocycle u_b(v/v(b))=u_a(v)/v(b), with bounded positive denominator. Completion is an all-semivaluation homeomorphism with exact primary contraction. Normalization branches are local normal centres above the marked point, not merely minimal components. Finite conductor value gives a singleton fibre; infinite value gives exactly the conductor quotient identification. The infinity stratum is closed and its finite complement open. Both C and D descend exactly iff constant on all fibres, including mixedzero and partialbottom. The same class-level pullback/descent result holds for every reduced finite injection over an excellent reduced local target. This supplies the cone rules asked in P12.2, not all geometric operators of P12.6.

Full BFJ specialization is proved on ALL semivaluations at a smooth affine closed germ over an algebraically closed field of characteristic zero. C is exactly its finite continuous subclass; D is its full canonically retracted formal psh class with identically bottom adjoined. Smooth DVR and empty dimension zero are direct. Source domains are not enlarged to positive-characteristic singular SNC settings.

The cone is not identified with literal classical or CLD logarithmic radial extensions. The mixedzero nodal profile (-1,0) belongs to this boundary cone and remains excluded by the stated radial compatibility obstruction. Global projective trivially valued energy constructions are not claimed to solve this local target.

P12.3 obstacle envelopes, P12.4 bounded mixed Monge--Ampere products, P12.5 energy and P12.6 geometric envelope/product/operator functoriality are separate. They are not constructed here and their omission is not a P12.2 gap. This manuscript is a complete-solution candidate for ORIGINAL P12.2 only, requiring subsequent independent whole-original review.
\end{Verbatim}
\end{document}
